\documentclass[UTF8, 12pt, fontset=fandol, oneside]{ctexart}
\usepackage{researchnotes}

\title{第 5 章\quad 多项式}
\author{}
\date{}

\begin{document}

\maketitle

\section*{第 5 章\quad 多项式}

\subsection*{§5.1\quad 基本概念}

\subsubsection*{5.1.1\quad 一元多项式代数}

\noindent\textbf{1.\quad 定义}

设 $\mathbb{F}$ 是一个数域，$x$ 是未定元，$a_0,a_1,\cdots,a_n\in\mathbb{F}(n\ge 0,a_n\ne0)$，称形式表达式
\[
a_nx^n+a_{n-1}x^{n-1}+\cdots+a_1x+a_0
\]
为数域 $\mathbb{F}$ 上关于未定元 $x$ 的 $n$ 次多项式. 数域 $\mathbb{F}$ 上的一元多项式全体组成的集合记为 $\mathbb{F}[x]$.

\noindent\textbf{2.\quad 运算及运算法则}

(1) 加法：设 $f(x),g(x)$ 是 $\mathbb{F}$ 上两个多项式，适当添上若干个零，可设
\[
f(x)=a_nx^n+a_{n-1}x^{n-1}+\cdots+a_1x+a_0,
\]
\[
g(x)=b_nx^n+b_{n-1}x^{n-1}+\cdots+b_1x+b_0,
\]
定义 $f(x)$ 和 $g(x)$ 的加法如下：
\[
f(x)+g(x)=(a_n+b_n)x^n+(a_{n-1}+b_{n-1})x^{n-1}
+\cdots+(a_1+b_1)x+(a_0+b_0).
\]

(2) 数乘：设
\[
f(x)=a_nx^n+a_{n-1}x^{n-1}+\cdots+a_1x+a_0,
\]
又 $k$ 是 $\mathbb{F}$ 中的数，则定义 $k$ 和 $f(x)$ 的数乘为
\[
kf(x)=ka_nx^n+ka_{n-1}x^{n-1}+\cdots+ka_1x+ka_0.
\]

$\mathbb{F}[x]$ 在加法和数乘的定义下满足向量空间的 8 条公理，因此 $\mathbb{F}[x]$ 是 $\mathbb{F}$ 上的向量空间.

(3) 乘法：设 $f(x),g(x)\in\mathbb{F}[x]$，且
\[
f(x)=a_nx^n+a_{n-1}x^{n-1}+\cdots+a_1x+a_0,\quad a_n\ne0,
\]
\[
g(x)=b_mx^m+b_{m-1}x^{m-1}+\cdots+b_1x+b_0,\quad b_m\ne0,
\]
定义 $f(x)$ 和 $g(x)$ 的乘法为
\[
f(x)\cdot g(x)=c_{n+m}x^{n+m}+c_{n+m-1}x^{n+m-1}+\cdots+c_1x+c_0,
\]
其中
\[
c_{n+m}=a_nb_m,
\]
\[
c_{n+m-1}=a_{n-1}b_m+a_nb_{m-1},
\]
\[
\cdots\cdots\cdots
\]
\[
c_k=\sum_{i+j=k}a_ib_j=a_kb_0+a_{k-1}b_1+\cdots+a_1b_{k-1}+a_0b_k,
\]
\[
\cdots\cdots\cdots
\]
\[
c_0=a_0b_0.
\]

乘法适合下列法则：
\[
\begin{array}{ll}
\text{(i)} & f(x)g(x)=g(x)f(x);\\
\text{(ii)} & (f(x)g(x))h(x)=f(x)(g(x)h(x));\\
\text{(iii)} & (f(x)+g(x))h(x)=f(x)h(x)+g(x)h(x);\\
\text{(iv)} & k(f(x)g(x))=(kf(x))g(x)=f(x)(kg(x)).
\end{array}
\]

\noindent\textbf{3.\quad 多项式的次数及其性质}

设
\[
f(x)=a_nx^n+a_{n-1}x^{n-1}+\cdots+a_1x+a_0,\quad a_n\ne0,
\]
则定义 $f(x)$ 的次数为 $n$，记为 $\deg f(x)=n$. 约定零多项式的次数为 $-\infty$.

\noindent\textbf{性质}

(1) $\deg f(x)g(x)=\deg f(x)+\deg g(x)$;

(2) $\deg kf(x)=\deg f(x)\ (k\ne0)$;

(3) $\deg(f(x)+g(x))\le\max\{\deg f(x),\deg g(x)\}$.

\subsubsection*{5.1.2\quad 整除}

\noindent\textbf{1.\quad 定义}

设 $f(x),g(x)$ 是 $\mathbb{F}$ 上的多项式，若存在 $\mathbb{F}$ 上的多项式 $h(x)$，使得
\[
f(x)=g(x)h(x),
\]
则称 $g(x)$ 是 $f(x)$ 的因式，或称 $g(x)$ 可整除 $f(x)$（也称 $f(x)$ 可被 $g(x)$ 整除），记为 $g(x)\mid f(x)$.

\noindent\textbf{2.\quad 定理}

设 $f(x),g(x)\in\mathbb{F}[x]$，$g(x)\ne0$，则必存在唯一的 $q(x),r(x)\in\mathbb{F}[x]$，使得
\[
f(x)=g(x)q(x)+r(x),
\]
且 $\deg r(x)<\deg g(x)$.

\noindent\textbf{3.\quad 推论}

设 $f(x),g(x)\in\mathbb{F}[x]$，$g(x)\ne0$，则 $g(x)\mid f(x)$ 的充要条件是 $r(x)=0$.

\subsubsection*{5.1.3\quad 最大公因式}

\noindent\textbf{1.\quad 定义}

设 $f(x),g(x)$ 是 $\mathbb{F}$ 上的多项式，$d(x)$ 是 $f(x),g(x)$ 的公因式，若对 $f(x),g(x)$ 的任一公因式 $h(x)$，都有 $h(x)\mid d(x)$，则称 $d(x)$ 为 $f(x),g(x)$ 的最大公因式，记为 $d(x)=(f(x),g(x))$.

特别地，若 $d(x)=1$，则称 $f(x),g(x)$ 互素.

\noindent\textbf{2.\quad 定理}

设 $f(x),g(x)$ 是 $\mathbb{F}$ 上的多项式，$d(x)$ 是它们的最大公因式，则必存在 $\mathbb{F}$ 上的多项式 $u(x),v(x)$，使得 $f(x)u(x)+g(x)v(x)=d(x)$.

\noindent\textbf{3.\quad 定理}

设 $f(x),g(x)$ 是 $\mathbb{F}$ 上的多项式，则它们互素的充要条件是，存在 $\mathbb{F}$ 上的多项式 $u(x),v(x)$，使得 $f(x)u(x)+g(x)v(x)=1$.

\noindent\textbf{4.\quad 互素多项式的性质}

(1) 若 $f_1(x)\mid g(x)$，$f_2(x)\mid g(x)$，且 $(f_1(x),f_2(x))=1$，则 $f_1(x)f_2(x)\mid g(x)$.

(2) 若 $(f(x),g(x))=1$，且 $f(x)\mid g(x)h(x)$，则 $f(x)\mid h(x)$.

(3) 若 $(f(x),g(x))=d(x)$，$f(x)=f_1(x)d(x)$，$g(x)=g_1(x)d(x)$，则 $(f_1(x),g_1(x))=1$.

(4) 若 $(f(x),g(x))=d(x)$，则 $(t(x)f(x),t(x)g(x))=t(x)d(x)$.

(5) 若 $(f_1(x),g(x))=1$，$(f_2(x),g(x))=1$，则 $(f_1(x)f_2(x),g(x))=1$.

\subsubsection*{5.1.4\quad 因式分解}

\noindent\textbf{1.\quad 定义}

设 $f(x)$ 是数域 $\mathbb{F}$ 上的多项式，若 $f(x)$ 可以分解为两个次数小于 $f(x)$ 的 $\mathbb{F}$ 上多项式之积，则称 $f(x)$ 是 $\mathbb{F}$ 上的可约多项式，否则称 $f(x)$ 为 $\mathbb{F}$ 上的不可约多项式.

\noindent\textbf{2.\quad 定理}

设 $p(x)$ 是 $\mathbb{F}$ 上的不可约多项式且 $p(x)\mid f(x)g(x)$，则或者 $p(x)\mid f(x)$ 或者 $p(x)\mid g(x)$.

\noindent\textbf{3.\quad 定理}

设 $f(x)$ 是数域 $\mathbb{F}$ 上的多项式且次数大于等于 1，则

(1) $f(x)$ 可分解为 $\mathbb{F}$ 上有限个不可约多项式之积;

(2) 若
\[
f(x)=p_1(x)p_2(x)\cdots p_s(x)=q_1(x)q_2(x)\cdots q_t(x)
\]
是 $f(x)$ 的两个不可约分解，则 $s=t$ 且经过适当调换不可约因式的次序后，有
\[
p_i(x)=c_iq_i(x)\quad (1\le i\le s),
\]
其中 $c_i$ 是 $\mathbb{F}$ 中的非零常数.

在多项式 $f(x)$ 的不可约分解中，若要求每个不可约因式都是首一多项式，且相同的不可约因式合并在一起，则
\[
f(x)=cp_1(x)^{e_1}p_2(x)^{e_2}\cdots p_k(x)^{e_k},
\]
其中 $c$ 是 $\mathbb{F}$ 中的非零常数. 上述分解式称为 $f(x)$ 的标准分解.

\noindent\textbf{4.\quad 定理}

多项式 $f(x)$ 没有重因式的充要条件是 $f(x)$ 和它的形式导数 $f'(x)$ 互素.

\noindent\textbf{5.\quad 定理}

设 $(f(x),f'(x))=d(x)$，则 $f(x)/d(x)$ 是一个没有重因式的多项式，且这个多项式的不可约因式和 $f(x)$ 的不可约因式相同（不计重数）.

\subsubsection*{5.1.5\quad 多项式函数}

\noindent\textbf{1.\quad 定义}

若 $f(b)=0$，则称 $b\in\mathbb{F}$ 适合 $\mathbb{F}$ 上的多项式 $f(x)$，也称 $b$ 为多项式 $f(x)$ 的零点或根.

\noindent\textbf{2.\quad 余数定理}

设 $f(x)\in\mathbb{F}[x]$，$b\in\mathbb{F}$，则存在 $\mathbb{F}$ 上的多项式 $g(x)$，使得
\[
f(x)=(x-b)g(x)+f(b).
\]

特别地，$b$ 是 $f(x)$ 的根的充要条件是 $(x-b)\mid f(x)$.

\noindent\textbf{3.\quad 定理}

设 $f(x)$ 是数域 $\mathbb{F}$ 上的 $n$ 次多项式，则 $f(x)$ 在 $\mathbb{F}$ 中最多只有 $n$ 个根.

\subsubsection*{5.1.6\quad 复系数多项式}

\noindent\textbf{1.\quad 代数基本定理}

次数大于零的复数域上的一元多项式至少有一个复数根.

\noindent\textbf{2.\quad 推论}

复数域上的一元 $n(n>0)$ 次多项式恰有 $n$ 个复数根（包括重根）.

\noindent\textbf{3.\quad 推论}

复数域上的一元 $n(n>0)$ 次多项式必可分解为一次因式的乘积.

\noindent\textbf{4.\quad 推论}

复数域上的不可约多项式都是一次多项式.

\noindent\textbf{5.\quad Vieta 定理}

若数域 $\mathbb{F}$ 上的多项式
\[
f(x)=x^n+p_1x^{n-1}+\cdots+p_{n-1}x+p_n
\]
在 $\mathbb{F}$ 中有 $n$ 个根 $x_1,x_2,\cdots,x_n$，则
\[
\sum_{i=1}^n x_i=x_1+x_2+\cdots+x_n=-p_1,
\]
\[
\sum_{i<j}x_ix_j=x_1x_2+\cdots+x_1x_n+x_2x_3+\cdots+x_{n-1}x_n=p_2,
\]
\[
\cdots\cdots\cdots
\]
\[
x_1x_2\cdots x_n=(-1)^n p_n.
\]

\subsubsection*{5.1.7\quad 实系数多项式}

\noindent\textbf{1.\quad 定理}

设有实系数多项式
\[
f(x)=a_nx^n+a_{n-1}x^{n-1}+\cdots+a_1x+a_0,
\]
若虚数 $a+bi$ 是 $f(x)$ 的根，则其共轭虚数 $a-bi$ 也是 $f(x)$ 的根.

\noindent\textbf{2.\quad 推论}

实数域上的不可约多项式的次数不超过 2. 因此实数域上的多项式可分解为有限个次数不超过 2 的实系数不可约多项式之积.

\subsubsection*{5.1.8\quad 有理系数多项式}

\noindent\textbf{1.\quad 定理}

设有整系数多项式
\[
f(x)=a_nx^n+a_{n-1}x^{n-1}+\cdots+a_1x+a_0,
\]
则有理数 $\dfrac{p}{q}$（$p,q$ 互素）是 $f(x)$ 的根的必要条件是 $p\mid a_0,\ q\mid a_n$.

\noindent\textbf{2.\quad 定理}

若整系数多项式 $f(x)$ 在有理数域上可约，则它在整数环上也可约，即可分解为两个次数较低的整系数多项式之积.

\noindent\textbf{3.\quad Eisenstein 判别法}

设有整系数多项式
\[
f(x)=a_nx^n+a_{n-1}x^{n-1}+\cdots+a_1x+a_0,\quad a_n\ne0,\ n>0,
\]
若有素数 $p$ 适合 $p\mid a_i(0\le i\le n-1)$，但 $p$ 不能整除 $a_n$，$p^2$ 也不能整除 $a_0$，则 $f(x)$ 在有理数域上不可约.

\subsubsection*{5.1.9\quad 多元多项式}

\noindent\textbf{1.\quad 定理}

设 $f(x_1,x_2,\cdots,x_n)$ 是数域 $\mathbb{F}$ 上的 $n$ 元非零多项式，则必存在 $\mathbb{F}$ 上元素 $a_1,a_2,\cdots,a_n$，使得
\[
f(a_1,a_2,\cdots,a_n)\ne0.
\]

\noindent\textbf{2.\quad 定义}

设 $f(x_1,x_2,\cdots,x_n)$ 是数域 $\mathbb{F}$ 上的 $n$ 元多项式，若对任意的 $1\le i<j\le n$，均有
\[
f(x_1,\cdots,x_i,\cdots,x_j,\cdots,x_n)
=f(x_1,\cdots,x_j,\cdots,x_i,\cdots,x_n),
\]
则称 $f(x_1,x_2,\cdots,x_n)$ 是数域 $\mathbb{F}$ 上的 $n$ 元对称多项式.

\noindent\textbf{3.\quad 定义}

下列多项式称为 $n$ 元初等对称多项式：
\[
\sigma_1=x_1+x_2+\cdots+x_n=\sum_{i=1}^n x_i,
\]
\[
\sigma_2=x_1x_2+\cdots+x_1x_n+x_2x_3+\cdots+x_{n-1}x_n
=\sum_{1\le i<j\le n}x_ix_j,
\]
\[
\cdots\cdots\cdots
\]
\[
\sigma_n=x_1x_2\cdots x_n.
\]

\noindent\textbf{4.\quad 对称多项式基本定理}

设 $f(x_1,x_2,\cdots,x_n)$ 是数域 $\mathbb{F}$ 上的 $n$ 元对称多项式，则存在唯一的 $\mathbb{F}$ 上的多项式 $g(y_1,y_2,\cdots,y_n)$，使得
\[
f(x_1,x_2,\cdots,x_n)=g(\sigma_1,\sigma_2,\cdots,\sigma_n).
\]

\noindent\textbf{5.\quad Newton 公式}

令 $s_k=x_1^k+x_2^k+\cdots+x_n^k$，若 $k<n$，则
\[
s_k-s_{k-1}\sigma_1+s_{k-2}\sigma_2+\cdots+(-1)^{k-1}s_1\sigma_{k-1}+(-1)^k k\sigma_k=0;
\]
若 $k\ge n$，则
\[
s_k-s_{k-1}\sigma_1+s_{k-2}\sigma_2+\cdots+(-1)^n s_{k-n}\sigma_n=0.
\]

\subsubsection*{5.1.10\quad 结式与判别式}

\noindent\textbf{1.\quad 定义}

设
\[
\begin{aligned}
f(x)&=a_0x^n+a_1x^{n-1}+\cdots+a_{n-1}x+a_n,\\
g(x)&=b_0x^m+b_1x^{m-1}+\cdots+b_{m-1}x+b_m,
\end{aligned}
\]
定义下列 $m+n$ 阶行列式：
\[
R(f,g)=
\left|\begin{array}{ccccccccc}
a_0 & a_1 & a_2 & \cdots & \cdots & a_n & 0 & \cdots & 0\\
0 & a_0 & a_1 & \cdots & \cdots & a_{n-1} & a_n & \cdots & 0\\
0 & 0 & a_0 & \cdots & \cdots & a_{n-2} & a_{n-1} & \cdots & 0\\
\vdots & \vdots & \vdots & \vdots & \vdots & \vdots & \vdots & \vdots & \vdots\\
0 & 0 & \cdots & 0 & a_0 & \cdots & \cdots & \cdots & a_n\\
b_0 & b_1 & b_2 & \cdots & \cdots & \cdots & b_m & \cdots & 0\\
0 & b_0 & b_1 & \cdots & \cdots & \cdots & b_{m-1} & b_m & \cdots\\
\vdots & \vdots & \vdots & \vdots & \vdots & \vdots & \vdots & \vdots & \vdots\\
0 & \cdots & 0 & b_0 & b_1 & \cdots & \cdots & \cdots & b_m
\end{array}\right|
\]
为 $f(x)$ 和 $g(x)$ 的结式.

\noindent\textbf{2.\quad 定理}

多项式 $f(x),g(x)$ 在复数域中有公共根的充要条件是它们的结式 $R(f,g)=0$. 等价地，$f(x),g(x)$ 互素的充要条件是它们的结式 $R(f,g)\ne0$.

\noindent\textbf{3.\quad 结式的其他表达式}

设
\[
\begin{aligned}
f(x)&=a_0x^n+a_1x^{n-1}+\cdots+a_{n-1}x+a_n,\\
g(x)&=b_0x^m+b_1x^{m-1}+\cdots+b_{m-1}x+b_m,
\end{aligned}
\]
又设 $f(x)$ 的根为 $x_1,x_2,\cdots,x_n$，$g(x)$ 的根为 $y_1,y_2,\cdots,y_m$，则
\[
R(f,g)=a_0^m\prod_{i=1}^n g(x_i)
=a_0^m b_0^n\prod_{i=1}^n\prod_{j=1}^m(x_i-y_j).
\]

\noindent\textbf{4.\quad 定义}

设有多项式 $f(x)=a_0x^n+a_1x^{n-1}+\cdots+a_{n-1}x+a_n$，定义 $f(x)$ 的判别式为
\[
\Delta(f(x))=(-1)^{\frac12 n(n-1)}a_0^{-1}R(f,f').
\]

\noindent\textbf{5.\quad 判别式的其他表达式}

设 $f(x)=a_0x^n+a_1x^{n-1}+\cdots+a_{n-1}x+a_n$，则 $f(x)$ 的判别式为
\[
\Delta(f(x))=a_0^{2n-2}\prod_{1\le i<j\le n}(x_i-x_j)^2,
\]
其中 $x_1,x_2,\cdots,x_n$ 是 $f(x)$ 的根.

\noindent\textbf{6.\quad 定理}

多项式 $f(x)$ 有重根的充要条件是其判别式等于零.

\subsection*{§5.2\quad 整除和带余除法}

带余除法和待定系数法是讨论整除问题的两个基本方法，下面是两个简单的例子.

\noindent\textbf{例 5.1}\quad 设 $g(x)=ax+b\in\mathbb{F}[x]$ 且 $a\ne0$，又 $f(x)\in\mathbb{F}[x]$，求证：$g(x)\mid f(x)^2$ 的充要条件是 $g(x)\mid f(x)$.

\noindent\textbf{证明 1}\quad 充分性显然，只需证明必要性. 设 $f(x)=g(x)q(x)+r$，则
\[
f(x)^2=g(x)^2q(x)^2+2rg(x)q(x)+r^2.
\]
由 $g(x)\mid f(x)^2$ 可得 $g(x)\mid r^2$，故 $r^2=0$，即 $r=0$，从而 $g(x)\mid f(x)$.

\noindent\textbf{证明 2}\quad 由余数定理，$f\left(-\dfrac{b}{a}\right)^2=0$，故 $f\left(-\dfrac{b}{a}\right)=0$，从而 $g(x)\mid f(x)$. \hfill $\square$

\noindent\textbf{例 5.2}\quad 设 $g(x)=ax^2+bx+c(abc\ne0)$，$f(x)=x^3+px^2+qx+r$，满足 $g(x)\mid f(x)$，求证：
\[
\frac{ap-b}{a}=\frac{aq-c}{b}=\frac{ar}{c}.
\]

\noindent\textbf{证明}\quad 用待定系数法，设
\[
\begin{aligned}
x^3+px^2+qx+r
&=(ax^2+bx+c)(mx+n)\\
&=amx^3+(an+bm)x^2+(bn+cm)x+cn.
\end{aligned}
\]
比较系数得
\[
am=1,\quad an+bm=p,\quad bn+cm=q,\quad cn=r.
\]
由此即可得到所需等式. \hfill $\square$

``凑项法''是指在要证明的等式中添加若干项再减去若干项来证明结论的方法，下面是两个典型的例子，请读者仔细领会.

\noindent\textbf{例 5.3}\quad 证明：$x^d-a^d$ 整除 $x^n-a^n$ 的充要条件是 $d\mid n$，其中 $a\ne0$.

\noindent\textbf{证明}\quad 充分性显然，现在来证明必要性. 若 $n=dq+r$，$0<r<d$，则
\[
x^n-a^n=x^n-x^ra^{dq}+x^ra^{dq}-a^n
=x^r(x^{dq}-a^{dq})+a^{dq}(x^r-a^r).
\]
注意到 $x^{dq}-a^{dq}$ 可被 $x^d-a^d$ 整除，而 $x^r-a^r$ 不能被 $x^d-a^d$ 整除，故 $x^n-a^n$ 不能被 $x^d-a^d$ 整除. \hfill $\square$

\noindent\textbf{例 5.4}\quad 设 $f(x)=x^{3m}+x^{3n+1}+x^{3p+2}$，其中 $m,n,p$ 为自然数，又 $g(x)=x^2+x+1$，求证：$g(x)\mid f(x)$.

\noindent\textbf{证明}\quad 首先注意这样一个事实：对任意的自然数 $k$，$x^{3k}-1$ 含因子 $x^3-1$，因此 $x^{3k}-1$ 总能被 $x^2+x+1$ 整除. 考虑下列等式：
\[
x^{3m}+x^{3n+1}+x^{3p+2}
=(x^{3m}-1)+x(x^{3n}-1)+x^2(x^{3p}-1)+(x^2+x+1),
\]
即知结论成立. \hfill $\square$

\subsection*{§5.3\quad 最大公因式与互素多项式}

若 $f(x),g(x)$ 的最大公因式是 $d(x)$，则必存在多项式 $u(x),v(x)$，使得 $f(x)u(x)+g(x)v(x)=d(x)$. 但读者需注意，这不是 $d(x)$ 为 $f(x),g(x)$ 最大公因式的充要条件，下面的例 5.5 说明了这一点. 另外，$u(x)$ 与 $v(x)$ 也不唯一，只有在一定的条件下才能保证唯一性，请参考例 5.6. 例 5.7 及其推论是对两个多项式的结论的推广.

\noindent\textbf{例 5.5}\quad 设 $d(x)=f(x)u(x)+g(x)v(x)$，举例说明 $d(x)$ 不必是 $f(x)$ 和 $g(x)$ 的最大公因式. 若进一步有 $d(x)\mid f(x)$，$d(x)\mid g(x)$，求证：$d(x)$ 必是 $f(x)$ 和 $g(x)$ 的最大公因式.

\noindent\textbf{证明}\quad 举例非常简单，请读者自己完成. 如果同时 $d(x)\mid f(x)$，$d(x)\mid g(x)$，则 $d(x)$ 是 $f(x)$ 和 $g(x)$ 的公因式. 若 $h(x)$ 也是 $f(x),g(x)$ 的公因式，则由 $h(x)\mid f(x)$，$h(x)\mid g(x)$ 可推出 $h(x)\mid (f(x)u(x)+g(x)v(x))=d(x)$，因此 $d(x)$ 是最大公因式. \hfill $\square$

\noindent\textbf{例 5.6}\quad 设 $f(x),g(x)$ 是次数不小于 1 的互素多项式，求证：必唯一地存在两个多项式 $u(x),v(x)$，使得
\[
f(x)u(x)+g(x)v(x)=1,
\]
且 $\deg u(x)<\deg g(x)$，$\deg v(x)<\deg f(x)$.

\noindent\textbf{证明}\quad 先证存在性. 因为 $(f(x),g(x))=1$，故必存在 $h(x),k(x)$，使得
\[
f(x)h(x)+g(x)k(x)=1.
\]
若 $h(x)$ 的次数大于等于 $g(x)$ 的次数，由带余除法，有
\[
h(x)=g(x)q(x)+u(x),\quad \deg u(x)<\deg g(x).
\]
代入前式得
\[
f(x)(g(x)q(x)+u(x))+g(x)k(x)=1,
\]
即有
\[
f(x)u(x)+g(x)(f(x)q(x)+k(x))=1.
\]
令 $v(x)=f(x)q(x)+k(x)$，则 $\deg v(x)<\deg f(x)$，否则由比较次数可知，上式将不可能成立.

再证唯一性. 设另有 $u_1(x),v_1(x)$ 适合条件，即
\[
f(x)u_1(x)+g(x)v_1(x)=f(x)u(x)+g(x)v(x)=1,
\]
则
\[
f(x)(u(x)-u_1(x))=g(x)(v_1(x)-v(x)).
\]
因为 $g(x)$ 和 $f(x)$ 互素，上式表明 $g(x)\mid (u(x)-u_1(x))$. 但 $u(x)-u_1(x)$ 的次数小于 $g(x)$ 的次数，所以只可能 $u(x)-u_1(x)=0$，即 $u(x)=u_1(x)$. 这时 $v(x)=v_1(x)$. \hfill $\square$

\noindent\textbf{例 5.7}\quad 设 $d(x)$ 是 $f_1(x),f_2(x),\cdots,f_m(x)$ 的最大公因式，求证：必存在多项式 $g_1(x),g_2(x),\cdots,g_m(x)$，使得
\[
f_1(x)g_1(x)+f_2(x)g_2(x)+\cdots+f_m(x)g_m(x)=d(x).
\]

\noindent\textbf{证明}\quad 用数学归纳法. 对 $m=2$，结论已成立. 设结论对 $m-1$ 成立. 设 $h(x)$ 是 $f_1(x),f_2(x),\cdots,f_{m-1}(x)$ 的最大公因式，则有 $g_1(x),g_2(x),\cdots,g_{m-1}(x)$，使得
\[
f_1(x)g_1(x)+f_2(x)g_2(x)+\cdots+f_{m-1}(x)g_{m-1}(x)=h(x).
\]
容易验证 $d(x)$ 是 $h(x)$ 和 $f_m(x)$ 的最大公因式，故存在 $u(x),v(x)$，使得
\[
h(x)u(x)+f_m(x)v(x)=d(x).
\]
将 $h(x)$ 代入可得
\[
f_1(x)g_1(x)u(x)+f_2(x)g_2(x)u(x)+\cdots+f_{m-1}(x)g_{m-1}(x)u(x)+f_m(x)v(x)=d(x),
\]
即知结论成立. \hfill $\square$

\noindent\textbf{推论}\quad 数域 $\mathbb{F}$ 上的多项式 $f_1(x),f_2(x),\cdots,f_m(x)$ 互素的充要条件是存在 $\mathbb{F}$ 上的多项式 $g_1(x),g_2(x),\cdots,g_m(x)$，使得
\[
f_1(x)g_1(x)+f_2(x)g_2(x)+\cdots+f_m(x)g_m(x)=1.
\]

下面几个例题是应用互素多项式性质的例子.

\noindent\textbf{例 5.8}\quad 设 $f(x)$ 和 $g(x)$ 互素，求证：$f(x^m)$ 和 $g(x^m)$ 也互素，其中 $m$ 为任一正整数.

\noindent\textbf{证明}\quad 因为 $f(x)$ 和 $g(x)$ 互素，故存在多项式 $u(x),v(x)$，使得
\[
f(x)u(x)+g(x)v(x)=1,
\]
从而有
\[
f(x^m)u(x^m)+g(x^m)v(x^m)=1,
\]
于是 $f(x^m)$ 和 $g(x^m)$ 互素. \hfill $\square$

\noindent\textbf{例 5.9}\quad 设 $(f(x),g(x))=1$，求证：$(f(x)g(x),f(x)+g(x))=1$.

\noindent\textbf{证明}\quad 由已知条件不难看出，$(f(x),f(x)+g(x))=1$，$(g(x),f(x)+g(x))=1$. 再由 §§ 5.1.3 中互素多项式的性质 (5) 即得 $(f(x)g(x),f(x)+g(x))=1$. \hfill $\square$

\noindent\textbf{例 5.10}\quad 设 $f_1(x),\cdots,f_m(x),g_1(x),\cdots,g_n(x)$ 为多项式，且
\[
(f_i(x),g_j(x))=1,\quad 1\le i\le m,\quad 1\le j\le n,
\]
求证：
\[
(f_1(x)f_2(x)\cdots f_m(x),g_1(x)g_2(x)\cdots g_n(x))=1.
\]

\noindent\textbf{证明}\quad 利用 §§ 5.1.3 中互素多项式的性质 (5) 以及归纳法即得结论. \hfill $\square$

\noindent\textbf{例 5.11}\quad 求证：$(f(x),g(x))=1$ 的充要条件是对任意给定的正整数 $m,n$，$(f(x)^m,g(x)^n)=1$.

\noindent\textbf{证明}\quad 必要性由例 5.10 即得. 反过来，若 $d(x)\ne1$ 是 $f(x)$ 和 $g(x)$ 的公因式，则它也是 $f(x)^m$ 和 $g(x)^n$ 的公因式，因此 $f(x)^m$ 和 $g(x)^n$ 不可能互素. \hfill $\square$

\noindent\textbf{例 5.12（中国剩余定理）}\quad 设 $g_1(x),\cdots,g_n(x)$ 是两两互素的多项式，$r_1(x),\cdots,r_n(x)$ 是 $n$ 个多项式. 求证：存在多项式 $f(x)$，使得 $f(x)=g_i(x)q_i(x)+r_i(x)(1\le i\le n)$.

\noindent\textbf{证明}\quad 先证存在多项式 $f_i(x)$，使得对任意的 $i$，有
\[
f_i(x)=g_i(x)p_i(x)+1,\quad g_j(x)\mid f_i(x)\ (j\ne i).
\]
一旦得证，只需令 $f(x)=r_1(x)f_1(x)+\cdots+r_n(x)f_n(x)$ 即可. 现构造 $f_1(x)$ 如下. 因为 $g_1(x)$ 和 $g_j(x)(j\ne1)$ 互素，故存在 $u_j(x),v_j(x)$，使得
\[
g_1(x)u_j(x)+g_j(x)v_j(x)=1.
\]
令
\[
f_1(x)=g_2(x)v_2(x)\cdots g_n(x)v_n(x)
=(1-g_1(x)u_2(x))\cdots(1-g_1(x)u_n(x)),
\]
显然 $f_1(x)$ 符合要求. 同理可构造 $f_i(x)$. \hfill $\square$

\noindent\textbf{例 5.13}\quad 设 $(f(x),g(x))=d(x)$，$[f(x),g(x)]=h(x)$，求证：
\[
(f(x)^n,g(x)^n)=d(x)^n,\quad [f(x)^n,g(x)^n]=h(x)^n.
\]

\noindent\textbf{证明}\quad 不妨设 $f(x),g(x)$ 都是首一多项式，$f(x)=f_1(x)d(x)$，$g(x)=g_1(x)d(x)$，则 $(f_1(x),g_1(x))=1$，$h(x)=f_1(x)g_1(x)d(x)$. 由例 5.11 可知，$(f_1(x)^n,g_1(x)^n)=1$，从而
\[
(f(x)^n,g(x)^n)=(f_1(x)^nd(x)^n,g_1(x)^nd(x)^n)=d(x)^n.
\]
从 $f(x)^ng(x)^n=(f(x)^n,g(x)^n)[f(x)^n,g(x)^n]=d(x)^n[f(x)^n,g(x)^n]$ 可得
\[
[f(x)^n,g(x)^n]=f_1(x)^ng_1(x)^nd(x)^n=h(x)^n.
\]
\hfill $\square$

\noindent\textbf{例 5.14}\quad 设 $f(x)=x^m-1$，$g(x)=x^n-1$，求证：$(f(x),g(x))=x^d-1$，其中 $d$ 是 $m,n$ 的最大公因子.

\noindent\textbf{证法 1}\quad 不妨设 $m\ge n$，$m=nq+r$，先证明 $(x^m-1,x^n-1)=(x^r-1,x^n-1)$. 假设 $d_1(x)=(x^m-1,x^n-1)$，$d_2(x)=(x^r-1,x^n-1)$. 注意到
\[
x^m-1=x^{nq+r}-1=x^r(x^{nq}-1)+(x^r-1),
\]
$(x^n-1)\mid (x^{nq}-1)$，故 $d_1(x)\mid (x^r-1)$，从而 $d_1(x)\mid d_2(x)$. 从上式也可以看出 $d_2(x)\mid (x^m-1)$，从而 $d_2(x)\mid d_1(x)$，因此 $d_1(x)=d_2(x)$. 又设 $n=qr+r_1$，则 $(x^m-1,x^n-1)=(x^n-1,x^r-1)=(x^r-1,x^{r_1}-1)$. 再由辗转相除，有某个 $r_{s-1}=q_{s+1}r_s$，其中 $r_s=d$ 是 $m,n$ 的最大公因子，于是 $(x^m-1,x^n-1)=(x^{r_{s-1}}-1,x^{r_s}-1)=x^d-1$.

\noindent\textbf{证法 2}\quad 只需求出 $f(x),g(x)$ 的公根. $f(x)$ 的根为
\[
\cos\frac{2k\pi}{m}+i\sin\frac{2k\pi}{m},\quad 1\le k\le m,
\]
$g(x)$ 的根为
\[
\cos\frac{2k\pi}{n}+i\sin\frac{2k\pi}{n},\quad 1\le k\le n,
\]
则公根为
\[
\cos\frac{2k\pi}{d}+i\sin\frac{2k\pi}{d},\quad 1\le k\le d.
\]
这就是 $x^d-1$ 的全部根，于是结论成立. \hfill $\square$

\noindent\textbf{例 5.15}\quad 设 $(f(x),g(x))=d(x)$，求证：对任意的正整数 $n$，
\[
(f(x)^n,f(x)^{n-1}g(x),\cdots,g(x)^n)=d(x)^n.
\]

\noindent\textbf{证明}\quad 显然 $d(x)^n$ 是 $f(x)^{n-k}g(x)^k(0\le k\le n)$ 的公因式. 又假设
\[
f(x)u(x)+g(x)v(x)=d(x),
\]
两边同时 $n$ 次方即可证明 $d(x)^n$ 是 $f(x)^{n-k}g(x)^k(0\le k\le n)$ 的最大公因式. \hfill $\square$

\subsection*{§5.4\quad 不可约多项式与因式分解}

数域 $\mathbb{F}$ 上的不可约多项式 $p(x)$ 与任一多项式 $f(x)$ 之间的关系很简单，即或者 $p(x)$ 整除 $f(x)$，或者 $p(x)$ 与 $f(x)$ 互素. 由不可约多项式的这一性质可以推出它的另一个重要性质，即例 5.16，通常称为不可约多项式的``素性''.

\noindent\textbf{例 5.16}\quad 设 $p(x)$ 是数域 $\mathbb{F}$ 上的非常数多项式，求证：$p(x)$ 为 $\mathbb{F}$ 上不可约多项式的充要条件是对 $\mathbb{F}$ 上任意适合 $p(x)\mid f(x)g(x)$ 的多项式 $f(x)$ 与 $g(x)$，或者 $p(x)\mid f(x)$，或者 $p(x)\mid g(x)$.

\noindent\textbf{证明}\quad 假设 $p(x)$ 不可约且 $p(x)\mid f(x)g(x)$. 若 $p(x)$ 不是 $f(x)$ 的因子，则 $p(x)$ 和 $f(x)$ 互素，因此 $p(x)\mid g(x)$.

反之，若 $p(x)$ 可约，设 $p(x)=p_1(x)p_2(x)$，其中 $\deg p_i(x)<\deg p(x)$，则 $p(x)\mid p_1(x)p_2(x)$，但 $p(x)$ 既不能整除 $p_1(x)$，又不能整除 $p_2(x)$. \hfill $\square$

\noindent\textbf{例 5.17}\quad 设 $f(x)$ 是数域 $\mathbb{F}$ 上的非常数多项式，求证：$f(x)$ 等于某个不可约多项式的幂的充要条件是对任意的非常数多项式 $g(x)$，或者 $f(x)$ 和 $g(x)$ 互素，或者 $f(x)$ 可以整除 $g(x)$ 的某个幂.

\noindent\textbf{证明}\quad 设 $f(x)=p(x)^k$，$p(x)$ 在 $\mathbb{F}$ 上不可约，且 $f(x)$ 和 $g(x)$ 不互素，则 $p(x)$ 是 $f(x)$ 和 $g(x)$ 的公因式，故 $f(x)$ 可以整除 $g(x)^k$.

反之，若 $f(x)=p(x)^m h(x)$，$p(x)$ 在 $\mathbb{F}$ 上不可约，$\deg h(x)>0$，且 $p(x)$ 不能整除 $h(x)$，则 $f(x)$ 既不和 $h(x)$ 互素，也不能整除 $h(x)$ 的任意次幂. \hfill $\square$

代数数是数论中的一个重要概念，与代数数密切相关的是它的极小多项式. 下面是有关代数数和极小多项式的最基本命题.

\noindent\textbf{例 5.18}\quad 设 $u$ 是复数域中某个数，若 $u$ 适合某个非零有理系数多项式（或整系数多项式）$f(x)=a_nx^n+a_{n-1}x^{n-1}+\cdots+a_1x+a_0$，则称 $u$ 是一个代数数. 证明：

(1) 对任一代数数 $u$，存在唯一一个 $u$ 适合的首一有理系数多项式 $g(x)$，使得 $g(x)$ 是 $u$ 适合的所有非零有理系数多项式中次数最小者. 这样的 $g(x)$ 称为 $u$ 的极小多项式或最小多项式.

(2) 设 $g(x)$ 是一个 $u$ 适合的首一有理系数多项式，则 $g(x)$ 是 $u$ 的极小多项式的充要条件是 $g(x)$ 是有理数域上的不可约多项式.

\noindent\textbf{证明}\quad (1) 在 $u$ 适合的所有非零有理系数多项式构成的集合中（由假设这个集合非空），可以取出一个次数最小的多项式，然后将其首一化，即可得到 $u$ 的极小多项式 $g(x)$. 为了证明极小多项式的唯一性，我们先证明极小多项式的一个基本性质，即极小多项式可以整除 $u$ 适合的任一多项式 $f(x)$. 假设
\[
f(x)=g(x)q(x)+r(x),\quad \deg r(x)<\deg g(x),
\]
则由 $f(u)=g(u)=0$ 可知 $r(u)=0$. 若 $r(x)\ne0$，则 $u$ 适合一个比 $g(x)$ 的次数更小的多项式 $r(x)$，这和 $g(x)$ 是极小多项式矛盾. 因此 $r(x)=0$，即 $g(x)\mid f(x)$. 设 $h(x)$ 也是 $u$ 的极小多项式，则由上述性质可得 $g(x)\mid h(x)$，$h(x)\mid g(x)$，从而 $g(x)$ 和 $h(x)$ 只差一个非零常数，又它们都是首一的，故只能相等，唯一性得证.

(2) 先证必要性. 若极小多项式 $g(x)$ 在有理数域上可约，则 $g(x)=g_1(x)g_2(x)$ 可分解为两个比 $g(x)$ 的次数更小的多项式的乘积. 由 $0=g(u)=g_1(u)g_2(u)$ 可知 $g_1(u)$ 和 $g_2(u)$ 中至少有一个等于零. 不妨设 $g_1(u)=0$，则 $u$ 适合一个比 $g(x)$ 的次数更小的多项式 $g_1(x)$，这和 $g(x)$ 是极小多项式矛盾. 再证充分性. 设 $g(x)$ 是 $u$ 适合的有理数域上的首一不可约多项式，$h(x)$ 是 $u$ 的极小多项式. 由极小多项式的基本性质可知 $h(x)\mid g(x)$，又 $g(x)$ 是不可约多项式，于是 $g(x)$ 和 $h(x)$ 只差一个非零常数，而它们又都是首一的，故只能相等. 因此 $g(x)$ 就是 $u$ 的极小多项式. \hfill $\square$

我们将在 §5.8 中给出代数数的极小多项式的例子及其应用. 对于一般域上的代数扩张（不限定是有理数域），在抽象代数课程中也会引入代数元的极小多项式的定义，其本质特征就是不可约性. 下面的例题告诉我们，数域 $\mathbb{F}$ 上的不可约多项式也满足例 5.18 中提及的极小多项式的基本性质.

\noindent\textbf{例 5.19}\quad 设 $p(x)$ 是数域 $\mathbb{F}$ 上的不可约多项式，$f(x)$ 是 $\mathbb{F}$ 上的多项式. 证明：若 $p(x)$ 的某个复根 $a$ 也是 $f(x)$ 的根，则 $p(x)\mid f(x)$. 特别地，$p(x)$ 的任一复根都是 $f(x)$ 的根.

\noindent\textbf{证明}\quad 若 $(p(x),f(x))=1$，则存在 $\mathbb{F}$ 上的多项式 $u(x),v(x)$，使得 $p(x)u(x)+f(x)v(x)=1$. 令 $x=a$ 可得 $1=p(a)u(a)+f(a)v(a)=0$，矛盾. 因此 $p(x)$ 与 $f(x)$ 不互素，从而只能是 $p(x)\mid f(x)$，结论得证. \hfill $\square$

多项式的标准分解是证明某些问题的有力工具，下面几个例子说明了这一点.

\noindent\textbf{例 5.20}\quad 证明：$g(x)^2\mid f(x)^2$ 的充要条件是 $g(x)\mid f(x)$.

\noindent\textbf{证明}\quad 充分性是显然的，只需证明必要性. 设 $f(x),g(x)$ 的公共标准分解为
\[
f(x)=cp_1(x)^{e_1}p_2(x)^{e_2}\cdots p_k(x)^{e_k},\quad
g(x)=dp_1(x)^{f_1}p_2(x)^{f_2}\cdots p_k(x)^{f_k},
\]
其中 $p_i(x)$ 为互不相同的首一不可约多项式，$c,d$ 是非零常数，则
\[
f(x)^2=c^2p_1(x)^{2e_1}p_2(x)^{2e_2}\cdots p_k(x)^{2e_k},\quad
g(x)^2=d^2p_1(x)^{2f_1}p_2(x)^{2f_2}\cdots p_k(x)^{2f_k}.
\]
若 $g(x)^2\mid f(x)^2$，则 $2f_i\le2e_i$，从而 $f_i\le e_i(1\le i\le k)$. 因此 $g(x)\mid f(x)$. \hfill $\square$

\noindent\textbf{例 5.21}\quad 设 $f(x),g(x)$ 是数域 $\mathbb{F}$ 上的多项式，若 $h(x)$ 是 $\mathbb{F}$ 上的多项式且适合 $f(x)\mid h(x)$，$g(x)\mid h(x)$，则称 $h(x)$ 是 $f(x)$ 和 $g(x)$ 的公倍式. 进一步，若对 $f(x)$ 和 $g(x)$ 的任一公倍式 $u(x)$，均有 $h(x)\mid u(x)$，则称 $h(x)$ 是 $f(x)$ 和 $g(x)$ 的最小公倍式，记为 $h(x)=[f(x),g(x)]$. 试证：
\[
f(x)g(x)=c(f(x),g(x))[f(x),g(x)],
\]
其中 $c$ 是 $\mathbb{F}$ 中的某个非零常数.

\noindent\textbf{证明}\quad 设 $f(x),g(x)$ 的公共标准分解为
\[
f(x)=c_1p_1(x)^{e_1}p_2(x)^{e_2}\cdots p_k(x)^{e_k},\quad
g(x)=c_2p_1(x)^{f_1}p_2(x)^{f_2}\cdots p_k(x)^{f_k},
\]
其中 $p_i(x)$ 为互不相同的首一不可约多项式，$c_1,c_2$ 是非零常数，则
\[
d(x)=p_1(x)^{r_1}p_2(x)^{r_2}\cdots p_k(x)^{r_k},\quad
h(x)=p_1(x)^{s_1}p_2(x)^{s_2}\cdots p_k(x)^{s_k},
\]
其中 $r_i=\min\{e_i,f_i\}$，$s_i=\max\{e_i,f_i\}$. 令 $c=c_1c_2$，则有
\[
f(x)g(x)=cd(x)h(x).
\]
\hfill $\square$

\noindent\textbf{例 5.13}\quad 设 $(f(x),g(x))=d(x)$，$[f(x),g(x)]=h(x)$，求证：
\[
(f(x)^n,g(x)^n)=d(x)^n,\quad [f(x)^n,g(x)^n]=h(x)^n.
\]

\noindent\textbf{证法 2}\quad 设 $f(x),g(x)$ 的公共标准分解为
\[
f(x)=cp_1(x)^{e_1}p_2(x)^{e_2}\cdots p_k(x)^{e_k},\quad
g(x)=dp_1(x)^{f_1}p_2(x)^{f_2}\cdots p_k(x)^{f_k},
\]
其中 $p_i(x)$ 为互不相同的首一不可约多项式，$c,d$ 是非零常数，则
\[
d(x)=p_1(x)^{r_1}p_2(x)^{r_2}\cdots p_k(x)^{r_k},\quad
h(x)=p_1(x)^{s_1}p_2(x)^{s_2}\cdots p_k(x)^{s_k},
\]
其中 $r_i=\min\{e_i,f_i\}$，$s_i=\max\{e_i,f_i\}$. 注意到
\[
f(x)^n=c^np_1(x)^{ne_1}p_2(x)^{ne_2}\cdots p_k(x)^{ne_k},\quad
g(x)^n=d^np_1(x)^{nf_1}p_2(x)^{nf_2}\cdots p_k(x)^{nf_k},
\]
并且 $\min\{ne_i,nf_i\}=nr_i$，$\max\{ne_i,nf_i\}=ns_i$，因此
\[
(f(x)^n,g(x)^n)=p_1(x)^{nr_1}p_2(x)^{nr_2}\cdots p_k(x)^{nr_k}=d(x)^n,
\]
\[
[f(x)^n,g(x)^n]=p_1(x)^{ns_1}p_2(x)^{ns_2}\cdots p_k(x)^{ns_k}=h(x)^n.
\]
\hfill $\square$

\subsection*{§5.5\quad 多项式函数与根}

判别多项式有无重根的最常用方法是判别该多项式及其形式导数是否有公因式. 下面是应用该方法的几个例子.

\noindent\textbf{例 5.22}\quad 证明：数域 $\mathbb{F}$ 上任意一个不可约多项式在复数域 $\mathbb{C}$ 中无重根.

\noindent\textbf{证明}\quad 设 $f(x)$ 是 $\mathbb{F}$ 上的不可约多项式，因为 $f(x)$ 不可能是 $f'(x)$ 的因式，故 $(f(x),f'(x))=1$，从而 $f(x)$ 在 $\mathbb{C}$ 中无重根. \hfill $\square$

\noindent\textbf{例 5.23}\quad 求证：$a$ 是多项式 $f(x)$ 的 $k$ 重根的充要条件是：
\[
f(a)=f'(a)=\cdots=f^{(k-1)}(a)=0,\quad f^{(k)}(a)\ne0.
\]

\noindent\textbf{证明}\quad 若 $a$ 是 $f(x)$ 的 $k$ 重根，可设 $f(x)=(x-a)^kg(x)$，$g(x)$ 不含因式 $x-a$. 通过对 $f(x)$ 求导可发现，$x-a$ 可整除 $f^{(j)}(x)(1\le j\le k-1)$. 因此
\[
f(a)=f'(a)=\cdots=f^{(k-1)}(a)=0.
\]
而 $f^{(k)}(a)=k!g(a)\ne0$，故必要性得证.

反之，若 $a$ 是 $f(x)$ 的 $m$ 重根，若 $m>k$，则由必要性的证明可知，将有 $f^{(k)}(a)=0$，这与已知矛盾. 同样，若 $m<k$，则由必要性的证明可知，将有 $f^{(m)}(a)\ne0$，这也与已知矛盾，于是只能 $m=k$. \hfill $\square$

\noindent\textbf{例 5.24}\quad 设 $\deg f(x)=n\ge1$，若 $f'(x)\mid f(x)$，证明：$f(x)$ 有 $n$ 重根.

\noindent\textbf{证法 1}\quad 设 $f(x)=\dfrac1n(x-a)f'(x)$，现证明 $a$ 是 $f(x)$ 的 $n$ 重根. 假设 $a$ 是 $f(x)$ 的 $k$ 重根，$f(x)=(x-a)^kg(x)$，$k<n$ 且 $g(x)$ 不含因式 $x-a$，则
\[
f'(x)=k(x-a)^{k-1}g(x)+(x-a)^kg'(x)=n(x-a)^{k-1}g(x).
\]
于是 $g(x)\mid (x-a)g'(x)$，而 $g(x)$ 与 $x-a$ 互素，故将有 $g(x)\mid g'(x)$. 引出矛盾.

\noindent\textbf{证法 2}\quad 设 $f(x)=\dfrac1n(x-a)f'(x)$，则
\[
\frac{f(x)}{(f(x),f'(x))}=b(x-a),\quad b\ne0.
\]
由 §§ 5.1.4 定理 5 可知，$x-a$ 是 $f(x)$ 仅有的不可约因式，因此 $f(x)=b(x-a)^n$. \hfill $\square$

一个一元 $n$ 次多项式在所在的数域内最多只有 $n$ 个根，利用这个命题可以证明一些有趣的结论，下面是 3 个例子.

\noindent\textbf{例 5.25}\quad 设 $f(x)$ 是数域 $\mathbb{F}$ 上的多项式，若对 $\mathbb{F}$ 中某个非零常数 $a$，有 $f(x+a)=f(x)$，求证：$f(x)$ 必是常数多项式.

\noindent\textbf{证明}\quad 假设 $f(x)$ 不是常数多项式，则 $f(x)-f(a)$ 也不是常数多项式，但由 $f(x+a)=f(x)$ 可知，$ka(k\in\mathbb{Z})$ 是 $f(x)-f(a)$ 的无穷多个根，矛盾. \hfill $\square$

\noindent\textbf{例 5.26}\quad 设 $f(x)$ 是非常数多项式且 $f(x)$ 可以整除 $f(x^m)(m>1)$，求证：$f(x)$ 的根只能是 $0$ 或 $1$ 的某个方根.

\noindent\textbf{证明}\quad 将 $f(x)$ 看成是复数域上的多项式，则 $f(x^m)=f(x)q(x)$. 假设 $c$ 是 $f(x)$ 的一个复根，即 $f(c)=0$，则 $f(c^m)=0$，即 $c^m$ 也是 $f(x)$ 的根. 由此可知 $c^{m^2},c^{m^3},\cdots$ 也都是 $f(x)$ 的根. 由于 $f(x)$ 只有有限个不同的复根，故存在正整数 $k,t$，使得 $c^{m^k}=c^{m^t}$. 因此若 $c\ne0$，则必存在某个正整数 $n$，使得 $c^n=1$. \hfill $\square$

\noindent\textbf{例 5.27}\quad 求证：$f(x)=\sin x$ 在实数域内不能表示为 $x$ 的多项式.

\noindent\textbf{证明}\quad 注意到 $f(x)=\sin x$ 在实数域内有无穷多个根，而任一非零多项式只能有有限个根，因此 $f(x)=\sin x$ 在实数域内不能表示为 $x$ 的多项式. \hfill $\square$

利用余数定理可以实现求根与判断整除性之间的相互转换，它常常使问题的解决变得简单. 下面是 3 个典型的例子.

\noindent\textbf{例 5.28}\quad 设 $n$ 是奇数，求证：$(x+y)(y+z)(x+z)$ 可整除 $(x+y+z)^n-x^n-y^n-z^n$.

\noindent\textbf{证明}\quad 将多项式 $(x+y+z)^n-x^n-y^n-z^n$ 看成是未定元 $x$ 的多项式. 当 $x=-y$ 时，$(x+y+z)^n-x^n-y^n-z^n=0$，因此 $x+y$ 是 $(x+y+z)^n-x^n-y^n-z^n$ 的因式. 同理 $x+z$，$y+z$ 也是因式. 又这 3 个因式互素，故 $(x+y)(y+z)(x+z)$ 可整除 $(x+y+z)^n-x^n-y^n-z^n$. \hfill $\square$

\noindent\textbf{例 5.29}\quad 设 $f(x)$ 是一个 $n$ 次多项式，若当 $k=0,1,\cdots,n$ 时有 $f(k)=\dfrac{k}{k+1}$，求 $f(n+1)$.

\noindent\textbf{解}\quad 令 $g(x)=(x+1)f(x)-x$，则 $0,1,\cdots,n$ 是 $g(x)$ 的根，因此
\[
g(x)=cx(x-1)(x-2)\cdots(x-n),
\]
即
\[
(x+1)f(x)-x=cx(x-1)(x-2)\cdots(x-n),
\]
其中 $c$ 是一个常数. 令 $x=-1$，可求出 $c=\dfrac{(-1)^{n+1}}{(n+1)!}$，从而
\[
f(x)=\frac{1}{x+1}\left(\frac{(-1)^{n+1}x(x-1)\cdots(x-n)}{(n+1)!}+x\right),
\]
故
\[
f(n+1)=\frac{1}{n+2}\left((-1)^{n+1}+n+1\right).
\]
当 $n$ 是奇数时，$f(n+1)=1$；当 $n$ 是偶数时，$f(n+1)=\dfrac{n}{n+2}$. \hfill $\square$

\noindent\textbf{例 5.30}\quad 设 $(x^4+x^3+x^2+x+1)\mid (x^3f_1(x^5)+x^2f_2(x^5)+xf_3(x^5)+f_4(x^5))$，这里 $f_i(x)(1\le i\le4)$ 都是实系数多项式，求证：$f_i(1)=0(1\le i\le4)$.

\noindent\textbf{证明}\quad 设 $\varepsilon_i(1\le i\le4)$ 是 1 的五次虚根，由条件可得
\[
\varepsilon_i^3f_1(1)+\varepsilon_i^2f_2(1)+\varepsilon_i f_3(1)+f_4(1)=0\quad (1\le i\le4).
\]
这是一个由 4 个未知数、4 个方程式组成的线性方程组（将 $f_i(1)$ 看成是未知数），其系数行列式是一个 Vandermonde 行列式，显然其值不等于零. 因此 $f_i(1)=0$. \hfill $\square$

\subsection*{§5.6\quad 复系数多项式}

Vieta 定理是多项式根与系数关系的最重要的结论，它有许多用途. 下面几个是应用 Vieta 定理的典型例子.

\noindent\textbf{例 5.31}\quad 设三次方程 $x^3+px^2+qx+r=0$ 的 3 个根成等差数列，求证：
\[
2p^3-9pq+27r=0.
\]

\noindent\textbf{证明}\quad 设方程的 3 个根为 $c-d,c,c+d$，则由 Vieta 定理可得
\[
\left\{
\begin{array}{l}
3c=-p,\\
3c^2-d^2=q,\\
c(c^2-d^2)=-r.
\end{array}
\right.
\]
由此可得 $2p^3-9pq+27r=0$. \hfill $\square$

\noindent\textbf{例 5.32}\quad 设三次方程 $x^3+px^2+qx+r=0(r\ne0)$ 的 3 个根成等比数列，求证：
\[
rp^3=q^3.
\]

\noindent\textbf{证明}\quad 设方程的 3 个根为 $\dfrac{c}{d},c,cd$，则由 Vieta 定理可得
\[
\left\{
\begin{array}{l}
\dfrac{c}{d}+c+cd=-p,\\
\dfrac{c^2}{d}+c^2+c^2d=q,\\
c^3=-r.
\end{array}
\right.
\]
由此可得 $rp^3=q^3$. \hfill $\square$

\noindent\textbf{例 5.33}\quad 设多项式 $x^3+3x^2+mx+n$ 的 3 个根成等差数列，多项式 $x^3-(m-2)x^2+(n-3)x+8$ 的 3 个根成等比数列，求 $m$ 和 $n$.

\noindent\textbf{解}\quad 由例 5.31 和例 5.32 可知，$m,n$ 应满足如下关系：
\[
\left\{
\begin{array}{l}
m=n+2,\\
-8(m-2)^3=(n-3)^3.
\end{array}
\right.
\]
若 $n-3=-2(m-2)$，则可联立求得 $m=3,n=1$.

若 $n-3=-2\omega(m-2)$，其中 $\omega=-\dfrac12+\dfrac{\sqrt3}{2}i$，则可联立求得 $m=2-\sqrt3 i$，$n=-\sqrt3 i$.

若 $n-3=-2\omega^2(m-2)$，则可联立求得 $m=2+\sqrt3 i$，$n=\sqrt3 i$. \hfill $\square$

\noindent\textbf{例 5.34}\quad 设 $x_1,x_2,x_3$ 是三次方程 $x^3+px^2+qx+r=0(r\ne0)$ 的 3 个根，求这 3 个根倒数的平方和.

\noindent\textbf{解}\quad 由 Vieta 定理可得
\[
\begin{aligned}
\frac{1}{x_1^2}+\frac{1}{x_2^2}+\frac{1}{x_3^2}
&=\frac{(x_1x_2+x_1x_3+x_2x_3)^2-2x_1x_2x_3(x_1+x_2+x_3)}
{x_1^2x_2^2x_3^2}\\
&=\frac{q^2-2pr}{r^2}.
\end{aligned}
\]
\hfill $\square$

\noindent\textbf{例 5.35}\quad 已知方程 $x^3+px^2+qx+r=0$ 的 3 个根为 $x_1,x_2,x_3$，求一个三次方程使其根为 $x_1^3,x_2^3,x_3^3$.

\noindent\textbf{解}\quad 由 Vieta 定理经计算可得
\[
\left\{
\begin{array}{l}
x_1^3+x_2^3+x_3^3=-p^3+3pq-3r,\\
x_1^3x_2^3+x_1^3x_3^3+x_2^3x_3^3=q^3-3pqr+3r^2,\\
x_1^3x_2^3x_3^3=-r^3.
\end{array}
\right.
\]
因此，以 $x_1^3,x_2^3,x_3^3$ 为根的三次方程为
\[
x^3+(p^3-3pq+3r)x^2+(q^3-3pqr+3r^2)x+r^3=0.
\]
\hfill $\square$

\noindent\textbf{例 5.36}\quad 设多项式 $x^3+px^2+qx+r$ 的 3 个根都是实数，求证：$p^2\ge3q$.

\noindent\textbf{证明}\quad 设多项式的 3 个根为 $x_1,x_2,x_3$，由已知条件可知：
\[
(x_1-x_2)^2+(x_2-x_3)^2+(x_1-x_3)^2\ge0.
\]
用 Vieta 定理可计算出
\[
(x_1-x_2)^2+(x_2-x_3)^2+(x_1-x_3)^2=2(p^2-3q).
\]
因此结论为真. \hfill $\square$

下面的例 5.37 也是讨论根与系数关系的，若用 Vieta 定理来做则比较繁琐，我们采用了不同的方法来处理.

\noindent\textbf{例 5.37}\quad 设 $f(x)=a_nx^n+a_{n-1}x^{n-1}+\cdots+a_1x+a_0$ 的 $n$ 个根 $x_1,x_2,\cdots,x_n$ 皆不等于零，求以 $\dfrac1{x_1},\dfrac1{x_2},\cdots,\dfrac1{x_n}$ 为根的多项式.

\noindent\textbf{解}\quad 令
\[
g(x)=a_0x^n+a_1x^{n-1}+\cdots+a_{n-1}x+a_n,
\]
则
\[
x_i^ng(x_i^{-1})=a_0+a_1x_i+\cdots+a_{n-1}x_i^{n-1}+a_nx_i^n=f(x_i)=0.
\]
因为 $x_i\ne0$，故 $g(x_i^{-1})=0$，即 $g(x)$ 的根为 $f(x)$ 根之倒数. \hfill $\square$

\noindent\textbf{例 5.38}\quad 设 $f(x)=a_nx^n+a_{n-1}x^{n-1}+\cdots+a_1x+a_0(a_na_0\ne0)$ 是数域 $\mathbb{F}$ 上的可约多项式，求证：多项式 $g(x)=a_0x^n+a_1x^{n-1}+\cdots+a_{n-1}x+a_n$ 在 $\mathbb{F}$ 上也可约.

\noindent\textbf{证明}\quad 设 $f(x)=p(x)q(x)$，其中 $\deg p(x)=m$，$\deg q(x)=n-m$，$0<m<n$，则
\[
g(x)=x^nf\left(\frac1x\right)
=x^np\left(\frac1x\right)q\left(\frac1x\right)
=\left(x^mp\left(\frac1x\right)\right)\left(x^{n-m}q\left(\frac1x\right)\right),
\]
因此 $g(x)$ 也可约. \hfill $\square$

\subsection*{§5.7\quad 实系数多项式}

下面 3 个例题是实系数多项式的基本性质.

\noindent\textbf{例 5.39}\quad 设 $f(x)$ 是复数域上的多项式，若对任意的实数 $c$，$f(c)$ 总是实数，求证：$f(x)$ 是实系数多项式.

\noindent\textbf{证明}\quad 设 $f(x)=a_nx^n+a_{n-1}x^{n-1}+\cdots+a_1x+a_0$，分别令 $x=0,1,2,\cdots,n$，得到一个以 $a_n,a_{n-1},\cdots,a_1,a_0$ 为未知数、由 $n+1$ 个方程式组成的实系数线性方程组. 该方程组的系数行列式是一个非零的 Vandermonde 行列式，故方程组必有唯一解，且解为实数. 因此 $f(x)$ 是实系数多项式. \hfill $\square$

\noindent\textbf{例 5.40}\quad 证明：奇数次实系数多项式必有实数根.

\noindent\textbf{证明}\quad 实系数多项式的虚根总是成对出现的，因此奇数次实系数多项式必有实数根. \hfill $\square$

\noindent\textbf{例 5.41}\quad 设 $f(x)=a_nx^n+a_{n-1}x^{n-1}+\cdots+a_1x+a_0$ 是实系数多项式，求证：

(1) 若 $a_i(0\le i\le n)$ 全是正数或全是负数，则 $f(x)$ 没有非负实根.

(2) 若 $(-1)^ia_i(0\le i\le n)$ 全是正数或全是负数，则 $f(x)$ 没有非正实根.

(3) 若 $a_n>0$ 且 $(-1)^{n-i}a_i>0(0\le i\le n-1)$，则 $f(x)$ 没有非正实根；若 $a_n>0$ 且 $(-1)^{n-i}a_i\ge0(0\le i\le n-1)$，则 $f(x)$ 没有负实根.

\noindent\textbf{证明}\quad (1) 若 $a_i$ 全是正数且 $f(x)$ 有非负实根 $c\ge0$，代入后可得
\[
f(c)=a_nc^n+a_{n-1}c^{n-1}+\cdots+a_1c+a_0\ge a_0>0,
\]
这和 $c$ 是根矛盾，因此 $f(x)$ 没有非负实根. 同理可证 $a_i$ 全是负数的情形.

(2) 和 (3) 同理可证. \hfill $\square$

\noindent\textbf{例 5.42}\quad 令 $\Delta=\dfrac{q^2}{4}+\dfrac{p^3}{27}$ 是实系数三次方程 $x^3+px+q=0$ 的判别式，求证：

(1) 若 $\Delta>0$，则方程有 1 个实根和 2 个共轭复根;

(2) 若 $\Delta=0$，则方程有 3 个实根，其中 2 个根相同;

(3) 若 $\Delta<0$，则方程有 3 个互不相等的实根.

\noindent\textbf{证明}\quad 注意到本题中的 $\Delta$ 和三次方程用结式定义的判别式相差一个负数（参考例 5.65），故由例 5.70 即得本题结论. 本题也可用 Cardano 公式直接证明. \hfill $\square$

\noindent\textbf{例 5.43}\quad 求证：实系数方程 $x^3+px^2+qx+r=0$ 的根的实部全是负数的充要条件是
\[
p>0,\quad r>0,\quad pq>r.
\]

\noindent\textbf{证明}\quad 先证必要性. 设原方程的 3 个根为 $x_1,x_2,x_3$，其中 $x_1$ 是实数根，$x_1<0$. 另假设 $x_2=a+bi$，$x_3=a-bi$，$a<0$，则
\[
p=-(x_1+x_2+x_3)=-(x_1+2a)>0,\quad
r=-x_1x_2x_3=-x_1(a^2+b^2)>0.
\]
\[
\begin{aligned}
pq-r
&=-(x_1+2a)(x_1x_2+x_1x_3+x_2x_3)+x_1(a^2+b^2)\\
&=-(x_1+2a)(2x_1a+a^2+b^2)+x_1(a^2+b^2)\\
&=-2a((x_1+a)^2+b^2)>0.
\end{aligned}
\]
又假设 $x_1,x_2,x_3$ 全是负实数，则显然 $p>0,q>0,r>0$，而
\[
\begin{aligned}
pq-r
&=-(x_1+x_2+x_3)(x_1x_2+x_1x_3+x_2x_3)+x_1x_2x_3\\
&=-(x_1^2+q)(x_2+x_3)>0.
\end{aligned}
\]

再证充分性. 由 $p>0,r>0,pq-r>0$ 可知 $q>0$，若方程的根是实数，则此根必是负数. 现假设方程有根 $x_1<0$，$x_2=a+bi$，$x_3=a-bi$，因为
\[
pq-r=-2a((x_1+a)^2+b^2)>0,
\]
故得 $a<0$，结论得证. \hfill $\square$

\noindent\textbf{例 5.44}\quad 设 $\varepsilon$ 是 1 的 $n$ 次根：
\[
\varepsilon=\cos\frac{2\pi}{n}+i\sin\frac{2\pi}{n},
\]
求证：$\varepsilon^{mi}(1\le i\le n)$ 是 $x^n-1=0$ 的全部根的充要条件是 $(m,n)=1$.

\noindent\textbf{证明}\quad 若 $(m,n)=1$，只要证明 $\varepsilon^{mi}(1\le i\le n)$ 互不相同即可. 若不然，有 $\varepsilon^{ms}=\varepsilon^{mt}(1\le s<t\le n)$，便有 $\varepsilon^{m(t-s)}=1$，$n\mid m(t-s)$. 因为 $n,m$ 互素，故 $n\mid (t-s)$，而 $0<t-s<n$，矛盾.

反之，若 $(m,n)=d>1$，则 $\varepsilon^{m\frac{n}{d}}=\varepsilon^{n\frac{m}{d}}=1$，从而 $\varepsilon^{mi}(1\le i\le n)$ 不可能是 $x^n-1=0$ 的全部根. \hfill $\square$

\noindent\textbf{例 5.45}\quad 设 $f(x)$ 是实系数首一多项式且无实数根，求证：$f(x)$ 可以表示为两个实系数多项式的平方和.

\noindent\textbf{证明}\quad 因为实系数多项式的虚根成对出现，故 $f(x)$ 是偶数次多项式，不妨设它的根为
\[
x_1,x_2,\cdots,x_n;\ \overline{x_1},\overline{x_2},\cdots,\overline{x_n}.
\]
令
\[
u(x)=(x-x_1)(x-x_2)\cdots(x-x_n),\quad
v(x)=(x-\overline{x_1})(x-\overline{x_2})\cdots(x-\overline{x_n}),
\]
则 $v(x)=\overline{u(x)}$，$f(x)=u(x)v(x)$. 又将 $u(x),v(x)$ 的实部和虚部分开，可设
\[
u(x)=g(x)+ih(x),\quad v(x)=g(x)-ih(x),
\]
即有
\[
f(x)=g(x)^2+h(x)^2.
\]
\hfill $\square$

\subsection*{§5.8\quad 有理系数多项式}

利用整数、有理数以及实数的性质来讨论有理系数多项式的性质是一种常用的方法，在下面的几个例子中读者将体会到这一点.

\noindent\textbf{例 5.46}\quad $f(x)$ 是次数大于零的首一整系数多项式，若 $f(0),f(1)$ 都是奇数，求证：$f(x)$ 没有有理根.

\noindent\textbf{证明}\quad 设 $f(x)=x^n+a_{n-1}x^{n-1}+\cdots+a_1x+a_0$. 从 $f(0)$ 是奇数可知 $a_0$ 是奇数，从 $f(1)$ 是奇数可知 $1+a_{n-1}+\cdots+a_1+a_0$ 是奇数. 假设 $f(x)$ 有有理根，因为 $f(x)$ 是首一的，故必有整数根，设其为 $c$，则
\[
c^n+a_{n-1}c^{n-1}+\cdots+a_1c+a_0=0.
\]
若 $c$ 是偶数，则上式左边为奇数，不可能等于零. 若 $c$ 是奇数，令 $c=2b+1$，其中 $b$ 是整数，可得
\[
(2b+1)^n+a_{n-1}(2b+1)^{n-1}+\cdots+a_1(2b+1)+a_0=0.
\]
用二项式定理展开后将看到，上式左边是一个偶数加上 $1+a_{n-1}+\cdots+a_1+a_0$，故必是奇数，也不可能等于零. 因此 $f(x)$ 没有有理根. \hfill $\square$

\noindent\textbf{例 5.47}\quad 设 $f(x)$ 是实系数多项式，若对任意的有理数 $c$，$f(c)$ 总是有理数，求证：$f(x)$ 是有理系数多项式.

\noindent\textbf{证明}\quad 类似例 5.39 可证明. \hfill $\square$

\noindent\textbf{例 5.48}\quad 设 $f(x)$ 是有理系数多项式，$a,b,c$ 是有理数，但 $\sqrt c$ 是无理数. 求证：若 $a+b\sqrt c$ 是 $f(x)$ 的根，则 $a-b\sqrt c$ 也是 $f(x)$ 的根.

\noindent\textbf{证明}\quad 设 $f(x)=a_nx^n+a_{n-1}x^{n-1}+\cdots+a_1x+a_0$，则
\[
f(a+b\sqrt c)=a_n(a+b\sqrt c)^n+a_{n-1}(a+b\sqrt c)^{n-1}+\cdots+a_1(a+b\sqrt c)+a_0=0.
\]
将 $(a+b\sqrt c)^k$ 用二项式定理展开，可设
\[
f(a+b\sqrt c)=A+B\sqrt c=0,
\]
其中 $A,B$ 都是有理数. 因为 $\sqrt c$ 是无理数，故 $A=B=0$. 因此
\[
f(a-b\sqrt c)=A-B\sqrt c=0,
\]
即 $a-b\sqrt c$ 也是 $f(x)$ 的根. \hfill $\square$

下面的例题和例 5.48 相似，但我们采用不同的方法来证明，这就是极小多项式的方法. 例 5.48 也可以用极小多项式的方法来证明，请读者自行思考.

\noindent\textbf{例 5.49}\quad 设 $f(x)$ 是有理系数多项式，$a,b,c,d$ 是有理数，但 $\sqrt c,\sqrt d,\sqrt{cd}$ 都是无理数. 求证：若 $a\sqrt c+b\sqrt d$ 是 $f(x)$ 的根，则下列数也是 $f(x)$ 的根：
\[
a\sqrt c-b\sqrt d,\quad -a\sqrt c+b\sqrt d,\quad -a\sqrt c-b\sqrt d.
\]

\noindent\textbf{证明}\quad 令
\[
g(x)=(x-(a\sqrt c+b\sqrt d))(x-(a\sqrt c-b\sqrt d))(x-(-a\sqrt c+b\sqrt d))(x-(-a\sqrt c-b\sqrt d)),
\]
则经计算可得
\[
g(x)=x^4-2(a^2c+b^2d)x^2+(a^2c-b^2d)^2.
\]
注意到 $g(x)$ 是一个有理系数首一多项式，只要证明它不可约，便可由例 5.18 得到 $g(x)$ 是 $a\sqrt c+b\sqrt d$ 的极小多项式，从而 $g(x)\mid f(x)$，于是结论成立. 显然 $g(x)$ 没有有理系数的一次因式，只要证明它没有有理系数的二次因式即可. 经过简单的计算可知，在 $g(x)$ 的 4 个一次因式中任取 2 个一次因式相乘都不是有理系数多项式，因此 $g(x)$ 没有有理系数的二次因式. \hfill $\square$

\noindent\textbf{例 5.50}\quad 求以 $\sqrt2+\sqrt[3]{3}$ 为根的次数最小的首一有理系数多项式.

\noindent\textbf{解}\quad 本题即求 $\sqrt2+\sqrt[3]{3}$ 的极小多项式. 令 $x-\sqrt2=\sqrt[3]{3}$，两边立方得到 $(x-\sqrt2)^3=3$. 整理可得 $x^3+6x-3=(3x^2+2)\sqrt2$，再两边平方可得，$\sqrt2+\sqrt[3]{3}$ 适合下列多项式：
\[
f(x)=x^6-6x^4-6x^3+12x^2-36x+1.
\]
由 $f(x)$ 的构造过程，不难看出 $f(x)$ 的 6 个根分别为 $\pm\sqrt2+\sqrt[3]{3}$，$\pm\sqrt2+\sqrt[3]{3}\omega$，$\pm\sqrt2+\sqrt[3]{3}\omega^2$，其中 $\omega=-\dfrac12+\dfrac{\sqrt3}{2}i$. 因此，我们有
\[
\begin{aligned}
f(x)={}&(x-\sqrt2-\sqrt[3]{3})(x+\sqrt2-\sqrt[3]{3})(x-\sqrt2-\sqrt[3]{3}\omega)(x+\sqrt2-\sqrt[3]{3}\omega)\\
&\cdot (x-\sqrt2-\sqrt[3]{3}\omega^2)(x+\sqrt2-\sqrt[3]{3}\omega^2).
\end{aligned}
\]
通过简单的验证可知，任取 $f(x)$ 的 2 个一次因式相乘都不是有理系数多项式；任取 $f(x)$ 的 3 个一次因式相乘也都不是有理系数多项式，因此 $f(x)$ 是有理数域上的不可约多项式，从而是 $\sqrt2+\sqrt[3]{3}$ 的极小多项式. \hfill $\square$

\noindent\textbf{例 5.51}\quad 求证：有理系数多项式 $x^4+px^2+q$ 在有理数域上可约的充要条件是或者 $p^2-4q=k^2$，其中 $k$ 是一个有理数；或者 $q$ 是某个有理数的平方，且 $\pm2\sqrt q-p$ 也是有理数的平方.

\noindent\textbf{证明}\quad 必要性. 若多项式 $x^4+px^2+q$ 在有理数域上可约，考虑下列两种情况：

(1) $x^4+px^2+q$ 有有理数根 $t$，这时 $t^2$ 是 $x^2+px+q$ 的有理根，因此其判别式 $p^2-4q$ 必是一个有理数的完全平方.

(2) $x^4+px^2+q$ 在有理数域上可分解为两个二次多项式的积. 设
\[
x^4+px^2+q=(x^2+ax+b)(x^2+cx+d),
\]
展开后比较系数可得
\[
\left\{
\begin{array}{l}
a+c=0,\\
ad+bc=0.
\end{array}
\right.
\]
若 $a=0$，则 $c=0$，这时将有 $p=b+d,q=bd$，因此 $p^2-4q=(b-d)^2$. 若 $a\ne0$，则 $b=d$，比较系数后可知 $p=2b-a^2$，$q=b^2$，因此 $\pm2\sqrt q-p=a^2$.

充分性. 若 $p^2-4q=k^2$，则
\[
x^4+px^2+q=x^4+px^2+\frac14(p+k)(p-k)
=\left(x^2+\frac12(p+k)\right)\left(x^2+\frac12(p-k)\right).
\]
因此原多项式可约.

若 $q=b^2$，$\pm2\sqrt q-p=\pm2b-p=a^2$，则 $p=-a^2\pm2b$. 于是
\[
x^4+px^2+q=x^4+(-a^2\pm2b)x^2+b^2=(x^2\pm b)^2-a^2x^2
\]
也可约. \hfill $\square$

Eisenstein 判别法是判定一个有理系数（或整系数）多项式不可约的最简单的方法之一. 如果不能直接用该判别法，可作一个变换（参考例 5.53）.

\noindent\textbf{例 5.52}\quad 设 $p_1,\cdots,p_m$ 是 $m$ 个互不相同的素数，求证：对任意的 $n\ge1$，下列多项式在有理数域上不可约：
\[
f(x)=x^n-p_1\cdots p_m.
\]

\noindent\textbf{证明}\quad 用 Eisenstein 判别法即可证明. \hfill $\square$

\noindent\textbf{例 5.53}\quad 证明：$x^8+1$ 在有理数域上不可约.

\noindent\textbf{证明}\quad 作代换 $x=y+1$，得
\[
x^8+1=(y+1)^8+1=y^8+8y^7+28y^6+56y^5+70y^4+56y^3+28y^2+8y+2.
\]
显然 2 可整除除第一项外的所有系数，但 4 不能整除常数项. 用 Eisenstein 判别法可知 $(y+1)^8+1$ 不可约，故 $x^8+1$ 也不可约. \hfill $\square$

\noindent\textbf{例 5.54}\quad 设 $f(x)$ 是有理系数多项式，已知 $\sqrt[n]{2}$ 是 $f(x)$ 的根，证明：$\sqrt[n]{2}\varepsilon,\sqrt[n]{2}\varepsilon^2,\cdots,\sqrt[n]{2}\varepsilon^{n-1}$ 也是 $f(x)$ 的根，其中 $\varepsilon=\cos\dfrac{2\pi}{n}+i\sin\dfrac{2\pi}{n}$ 是 1 的 $n$ 次根.

\noindent\textbf{证明}\quad 显然 $\sqrt[n]{2}$ 适合多项式 $x^n-2$，由 Eisenstein 判别法可知，$x^n-2$ 在有理数域上不可约，因此它是 $\sqrt[n]{2}$ 的极小多项式. 最后由极小多项式的基本性质可得 $(x^n-2)\mid f(x)$，从而结论得证. \hfill $\square$

\noindent\textbf{例 5.55}\quad 设 $f(x)$ 是次数大于 1 的奇数次有理系数不可约多项式，求证：若 $x_1,x_2$ 是 $f(x)$ 在复数域内两个不同的根，则 $x_1+x_2$ 必不是有理数.

\noindent\textbf{证明}\quad 不妨设 $f(x)$ 为首一多项式，我们用反证法来证明结论. 设 $x_1+x_2=r$ 为有理数，则有理系数多项式 $f(x)$ 与 $f(r-x)$ 有公共根 $x_1$. 因为 $f(x)$ 在有理数域上不可约，故 $f(x)$ 是 $x_1$ 的极小多项式，从而由极小多项式的基本性质可得 $f(x)\mid f(r-x)$. 注意到 $f(x)$ 与 $f(r-x)$ 次数相同，首项系数相反，从而有 $f(r-x)=-f(x)$. 令 $x=\dfrac r2$，则可得 $f\left(\dfrac r2\right)=0$，即 $\dfrac r2$ 是 $f(x)$ 的一个有理根，这与 $f(x)$ 在有理数域上不可约相矛盾. \hfill $\square$

除了用 Eisenstein 判别法外，另一个常用的方法是反证法，请参考下面两个例子.

\noindent\textbf{例 5.56}\quad 设 $f(x)=(x-a_1)(x-a_2)\cdots(x-a_n)-1$，其中 $a_1,a_2,\cdots,a_n$ 是 $n$ 个不同的整数，求证：$f(x)$ 在有理数域上不可约.

\noindent\textbf{证明}\quad 只要证明 $f(x)$ 在整数环上不可约即可. 用反证法，设 $f(x)=g(x)h(x)$，其中 $g(x),h(x)$ 都是次数小于 $n$ 的首一整系数多项式. 注意到
\[
g(a_i)h(a_i)=-1,
\]
因为 $g(x),h(x)$ 是整系数多项式，故 $g(a_i)=1,h(a_i)=-1$ 或 $g(a_i)=-1,h(a_i)=1$. 无论是哪种情况，都有
\[
g(a_i)+h(a_i)=0,\quad 1\le i\le n,
\]
即次数小于 $n$ 的多项式 $g(x)+h(x)$ 有 $n$ 个不同的根，故 $g(x)+h(x)=0$. 因此 $f(x)=-g(x)^2$，但 $f(x)$ 是首一多项式，而 $-g(x)^2$ 的首项系数为 $-1$，矛盾. \hfill $\square$

\noindent\textbf{例 5.57}\quad 设 $f(x)=(x-a_1)^2(x-a_2)^2\cdots(x-a_n)^2+1$，其中 $a_1,a_2,\cdots,a_n$ 是 $n$ 个不同的整数，求证：$f(x)$ 在有理数域上不可约.

\noindent\textbf{证明}\quad 只要证明 $f(x)$ 在整数环上不可约即可. 用反证法，设 $f(x)=u(x)v(x)$，其中 $u(x),v(x)$ 都是次数小于 $2n$ 的首一整系数多项式. 注意到 $f(x)$ 没有实根，故 $u(x),v(x)$ 也都没有实根，从而由实系数多项式虚根成对可知，$u(x),v(x)$ 作为实数域上的函数都恒大于零. 由于 $f(x)$ 是 $2n$ 次多项式，故 $u(x)$ 和 $v(x)$ 的次数至少有一个不超过 $n$，不妨设 $u(x)$ 的次数不超过 $n$.

若 $u(x)$ 的次数小于 $n$，则由 $f(a_i)=1$ 可得 $u(a_i)v(a_i)=1$，因此 $u(a_i)=1$. 考虑非零多项式 $u(x)-1$，由上面的分析可知它有 $n$ 个不同的根 $a_1,a_2,\cdots,a_n$，这与它的次数小于 $n$ 矛盾.

因此 $u(x)$ 只能是 $n$ 次首一多项式，于是 $v(x)$ 也是 $n$ 次首一多项式. 另一方面，由于 $u(a_i)v(a_i)=1$，故 $u(a_i)=v(a_i)=1(1\le i\le n)$. 注意到 $u(x)-v(x)$ 的次数小于 $n$，并且它有 $n$ 个不同的根 $a_1,a_2,\cdots,a_n$，因此只能是 $u(x)=v(x)$，$f(x)=u(x)^2$. 令 $h(x)=(x-a_1)(x-a_2)\cdots(x-a_n)$，则 $u(x)^2=h(x)^2+1$，即
\[
(u(x)+h(x))(u(x)-h(x))=1.
\]
因为 $u(x),h(x)$ 都是整系数多项式，故或者 $u(x)+h(x)=1,u(x)-h(x)=1$；或者 $u(x)+h(x)=-1,u(x)-h(x)=-1$，于是 $h(x)=0$，矛盾. 因此结论得证. \hfill $\square$

\subsection*{§5.9\quad 多元多项式}

数域 $\mathbb{K}$ 上的多元多项式有 3 个重要的性质. 首先是整性，即两个非零多元多项式的乘积仍然是非零多元多项式. 其次是多元多项式的非零性等价于多元函数的非零性. 最后是因式分解定理对多元多项式仍然成立，不过它的证明将在抽象代数课程中给出. 下面我们先来看多元多项式整性的若干应用.

\noindent\textbf{例 5.58}\quad 设 $f(x_1,\cdots,x_n),g(x_1,\cdots,x_n)\ne0$ 是 $\mathbb{K}$ 上的多元多项式. 假设对一切使 $g(a_1,\cdots,a_n)\ne0$ 的 $a_1,\cdots,a_n\in\mathbb{K}$，均有 $f(a_1,\cdots,a_n)=0$，求证：
\[
f(x_1,\cdots,x_n)=0.
\]

\noindent\textbf{证明}\quad 用反证法，假设 $f(x_1,\cdots,x_n)\ne0$，则由多元多项式的整性可知
\[
h(x_1,\cdots,x_n)=f(x_1,\cdots,x_n)g(x_1,\cdots,x_n)\ne0,
\]
于是存在 $a_1,\cdots,a_n\in\mathbb{K}$，使得 $h(a_1,\cdots,a_n)\ne0$，从而 $f(a_1,\cdots,a_n)\ne0$ 并且 $g(a_1,\cdots,a_n)\ne0$，这与假设矛盾. \hfill $\square$

我们可以利用多元多项式的整性去解决文字行列式求值中的两个问题.

\noindent\textbf{例 5.59}\quad 设 $A(x_1,x_2,\cdots,x_m)=(a_{ij})$ 为 $n$ 阶方阵，其元素 $a_{ij}=a_{ij}(x_1,x_2,\cdots,x_m)$ 都是 $\mathbb{K}$ 上的多元多项式. 设 $g(x_1,x_2,\cdots,x_m),h_i(x_1,x_2,\cdots,x_m)\ne0(1\le i\le k)$ 都是 $\mathbb{K}$ 上的多元多项式，
\[
U=\{(a_1,a_2,\cdots,a_m)\in\mathbb{K}^m\mid h_i(a_1,a_2,\cdots,a_m)\ne0(1\le i\le k)\}.
\]
若对所有的 $(a_1,a_2,\cdots,a_m)\in U$，都成立
\[
|A(a_1,a_2,\cdots,a_m)|=g(a_1,a_2,\cdots,a_m),
\]
证明：
\[
|A(x_1,x_2,\cdots,x_m)|=g(x_1,x_2,\cdots,x_m).
\]

\noindent\textbf{证明}\quad 用反证法，设 $(|A|-g)(x_1,x_2,\cdots,x_m)\ne0$，则由多元多项式的整性可知
\[
(|A|-g)h_1\cdots h_k(x_1,x_2,\cdots,x_m)\ne0,
\]
于是存在 $a_1,a_2,\cdots,a_m\in\mathbb{K}$，使得 $(|A|-g)h_1\cdots h_k(a_1,a_2,\cdots,a_m)\ne0$，从而 $h_i(a_1,a_2,\cdots,a_m)\ne0(1\le i\le k)$，即 $(a_1,a_2,\cdots,a_m)\in U$，并且 $(|A|-g)(a_1,a_2,\cdots,a_m)\ne0$，即 $|A(a_1,a_2,\cdots,a_m)|\ne g(a_1,a_2,\cdots,a_m)$，这与假设矛盾. \hfill $\square$

例 5.59 告诉我们：在元素为多元多项式的文字行列式的求值过程中，在假设某些非零条件成立的情形下得到的结果，其实就是所求行列式的值. 因此，在求行列式的过程中，可以暂不考虑未定元取特殊值的情形，而把主要精力放在一般的情形进行计算即可. 例如，读者不难发现在例 1.4，例 1.5，例 1.14 和例 1.26 的求值过程中，某些参数等于零的讨论并不影响最终的答案.

利用多元多项式的整性还可以给出第 1 章行列式求根法的严格证明.

\noindent\textbf{例 5.60}\quad 设 $A(x_1,x_2,\cdots,x_m)=(a_{ij})$ 为 $n$ 阶方阵，其元素 $a_{ij}=a_{ij}(x_1,x_2,\cdots,x_m)$ 都是 $\mathbb{K}$ 上的多元多项式，于是 $|A|$ 也是 $\mathbb{K}$ 上的多元多项式. 若把 $x_1$ 看成主未定元，则可将 $|A|$ 整理成关于 $x_1$ 的一元多项式：
\begin{equation}
|A|=c_0(x_2,\cdots,x_m)x_1^d+c_1(x_2,\cdots,x_m)x_1^{d-1}+\cdots+c_d(x_2,\cdots,x_m),
\tag{5.1}
\end{equation}
其中 $c_0(x_2,\cdots,x_m)\ne0$，$d\ge1$ 为次数. 假设存在互异的多项式 $g_1(x_2,\cdots,x_m),\cdots,g_d(x_2,\cdots,x_m)$，使得当 $x_1=g_i(x_2,\cdots,x_m)(1\le i\le d)$ 时 $|A|=0$，证明：
\[
|A|=c_0(x_2,\cdots,x_m)\cdot (x_1-g_1(x_2,\cdots,x_m))\cdots (x_1-g_d(x_2,\cdots,x_m)).
\]

\noindent\textbf{证明}\quad 由假设 $0=c_0(x_2,\cdots,x_m)g_1^d+\cdots+c_{d-1}(x_2,\cdots,x_m)g_1+c_d(x_2,\cdots,x_m)$，故有
\[
\begin{aligned}
|A|
&=c_0(x_2,\cdots,x_m)(x_1^d-g_1^d)+\cdots+c_{d-1}(x_2,\cdots,x_m)(x_1-g_1)\\
&=(x_1-g_1)R_1(x_1,x_2,\cdots,x_m).
\end{aligned}
\]
在上式中令 $x_1=g_2(x_2,\cdots,x_m)$，则有 $0=(g_2-g_1)R_1(g_2,x_2,\cdots,x_m)$. 注意到 $g_2-g_1\ne0$，故由多元多项式的整性可得 $R_1(g_2,x_2,\cdots,x_m)=0$. 再由相同的讨论可知，$x_1-g_2$ 是 $R_1(x_1,x_2,\cdots,x_m)$ 的因式，从而
\[
|A|=(x_1-g_1)(x_1-g_2)R_2(x_1,x_2,\cdots,x_m).
\]
不断地这样做下去，可得
\[
|A|=(x_1-g_1)\cdots(x_1-g_d)R_d(x_1,x_2,\cdots,x_m),
\]
最后与 (5.1) 式比较 $x_1$ 的首项系数可得 $R_d(x_1,x_2,\cdots,x_m)=c_0(x_2,\cdots,x_m)$. \hfill $\square$

\noindent\textbf{注}\quad 若首项系数 $c_0(x_2,\cdots,x_m)$ 能直接求出，或者可以继续通过求根法进行分解的话，那么 $|A|$ 的值就能计算出来. 请读者参考 §1.8 求根法中的例题.

对称多项式是一类常见的多元多项式，有着广泛的用途. 求对称多项式的初等对称多项式表示通常有两种方法，一种是对称多项式基本定理证明中的构造方法（适合次数较高时，配合待定系数法使用），另一种是 Vieta 定理（适合次数较低时使用）. 以下例题用了这两种方法来求解，请读者自行比较.

\noindent\textbf{例 5.61}\quad 求证：$f(x)=x^3+px^2+qx+r$ 某一根的平方等于其他两根平方和的充要条件是：
\[
p^4(p^2-2q)=2(p^3-2pq+2r)^2.
\]

\noindent\textbf{证明}\quad 设 $f(x)$ 的 3 个根为 $x_1,x_2,x_3$，则某一根的平方等于其他两根平方和的充要条件是：
\[
(x_1^2+x_2^2-x_3^2)(x_1^2+x_3^2-x_2^2)(x_2^2+x_3^2-x_1^2)=0.
\]

\noindent\textbf{证法 1}\quad 由 Vieta 定理，有
\[
\begin{aligned}
&(x_1^2+x_2^2-x_3^2)(x_1^2+x_3^2-x_2^2)(x_2^2+x_3^2-x_1^2)\\
={}&(x_1^2+x_2^2+x_3^2-2x_3^2)(x_1^2+x_2^2+x_3^2-2x_2^2)(x_1^2+x_2^2+x_3^2-2x_1^2)\\
={}&(p^2-2q-2x_3^2)(p^2-2q-2x_2^2)(p^2-2q-2x_1^2)\\
={}&(p^2-2q)^3-2(x_1^2+x_2^2+x_3^2)(p^2-2q)^2\\
&+4(x_1^2x_2^2+x_1^2x_3^2+x_2^2x_3^2)(p^2-2q)-8x_1^2x_2^2x_3^2\\
={}&(p^2-2q)^3-2(p^2-2q)(p^2-2q)^2+4(q^2-2pr)(p^2-2q)-8r^2\\
={}&-(p^2-2q)^3+4(q^2-2pr)(p^2-2q)-8r^2\\
={}&-p^6+6p^4q-8p^3r-8p^2q^2+16pqr-8r^2\\
={}&p^4(p^2-2q)-2(p^3-2pq+2r)^2,
\end{aligned}
\]
由此即得结论.

\noindent\textbf{证法 2}\quad 设 $g(x_1,x_2,x_3)=(x_1^2+x_2^2-x_3^2)(x_1^2+x_3^2-x_2^2)(x_2^2+x_3^2-x_1^2)$，则 $g$ 是六次齐次多项式，字典排序的首项为 $-x_1^6$，指数组为 $(6,0,0)$，对应的首项为 $-\sigma_1^6$. 设其余的指数组为 $(k_1,k_2,k_3)$，则有 $k_1+k_2+k_3=6$，$6\ge k_1\ge k_2\ge k_3\ge0$. 因此，可能的指数组分别为 $(5,1,0)$，$(4,2,0)$，$(4,1,1)$，$(3,3,0)$，$(3,2,1)$，$(2,2,2)$，对应的首项分别为 $\sigma_1^4\sigma_2$，$\sigma_1^2\sigma_2^2$，$\sigma_1^3\sigma_3$，$\sigma_2^3$，$\sigma_1\sigma_2\sigma_3$，$\sigma_3^2$. 用待定系数法，设
\[
g(x_1,x_2,x_3)=-\sigma_1^6+c_1\sigma_1^4\sigma_2+c_2\sigma_1^2\sigma_2^2+c_3\sigma_1^3\sigma_3+c_4\sigma_2^3+c_5\sigma_1\sigma_2\sigma_3+c_6\sigma_3^2.
\]
取 $x_1,x_2,x_3$ 的一些特殊值，可得到关于 $c_i(1\le i\le6)$ 的线性方程组，不难解得
\[
c_1=6,\quad c_2=c_3=c_6=-8,\quad c_4=0,\quad c_5=16.
\]
因此
\[
g(x_1,x_2,x_3)=-\sigma_1^6+6\sigma_1^4\sigma_2-8\sigma_1^2\sigma_2^2-8\sigma_1^3\sigma_3+16\sigma_1\sigma_2\sigma_3-8\sigma_3^2,
\]
经简单计算即得所求结论. \hfill $\square$

\noindent\textbf{注}\quad 类似例 5.61 的证明可知，例 5.31 和例 5.32 的结论其实是充要条件，请读者自行验证.

\noindent\textbf{例 5.62}\quad 设 $f(x_1,x_2,\cdots,x_n),g(x_1,x_2,\cdots,x_n)\ne0$ 为数域 $\mathbb{K}$ 上的多元多项式，称分式
\[
Q(x_1,x_2,\cdots,x_n)=\frac{f(x_1,x_2,\cdots,x_n)}{g(x_1,x_2,\cdots,x_n)}
\]
为关于未定元 $x_1,x_2,\cdots,x_n$ 的有理函数. 若对换任意两个未定元的位置，得到的有理函数保持不变，则称之为对称有理函数，例如 $\sum_{i=1}^n\dfrac1{x_i}$. 证明：对称有理函数可表示为初等对称多项式的有理函数.

\noindent\textbf{证明}\quad 设
\[
Q(x_1,x_2,\cdots,x_n)=\frac{f(x_1,x_2,\cdots,x_n)}{g(x_1,x_2,\cdots,x_n)}
\]
为对称有理函数，其中 $g(x_1,x_2,\cdots,x_n)\ne0$，于是对任意的全排列 $(k_1,k_2,\cdots,k_n)\in S_n$，$g(x_{k_1},x_{k_2},\cdots,x_{k_n})\ne0$，从而
\[
Q(x_1,x_2,\cdots,x_n)=
\frac{
f(x_1,x_2,\cdots,x_n)\displaystyle\prod_{(k_1,k_2,\cdots,k_n)\in S_n\setminus\{(1,2,\cdots,n)\}}g(x_{k_1},x_{k_2},\cdots,x_{k_n})
}{
\displaystyle\prod_{(k_1,k_2,\cdots,k_n)\in S_n}g(x_{k_1},x_{k_2},\cdots,x_{k_n})
}.
\]
注意到上述分式的分母 $\displaystyle\prod_{(k_1,k_2,\cdots,k_n)\in S_n}g(x_{k_1},x_{k_2},\cdots,x_{k_n})$ 是对称多项式，分子等于 $Q(x_1,x_2,\cdots,x_n)\displaystyle\prod_{(k_1,k_2,\cdots,k_n)\in S_n}g(x_{k_1},x_{k_2},\cdots,x_{k_n})$ 也是对称多项式，故由对称多项式基本定理可知结论成立. \hfill $\square$

Newton 公式和例 5.64 给出了幂和 $s_k$ 与初等对称多项式 $\sigma_k$ 之间的关系式，它们为一些求根问题的讨论提供了工具，比如例 5.63 以及第 6 章特征值的求解等.

\noindent\textbf{例 5.63}\quad 求解下列方程组：
\[
\left\{
\begin{array}{l}
x_1+x_2+x_3+x_4=4,\\
x_1^2+x_2^2+x_3^2+x_4^2=4,\\
x_1^3+x_2^3+x_3^3+x_4^3=4,\\
x_1^4+x_2^4+x_3^4+x_4^4=4.
\end{array}
\right.
\]

\noindent\textbf{解}\quad 上述方程组表明：$s_1=s_2=s_3=s_4=4$. 由 Newton 公式可求得 $\sigma_1=4$，$\sigma_2=6$，$\sigma_3=4$，$\sigma_4=1$. 因此若将 $x_1,x_2,x_3,x_4$ 看成是某个四次方程的根，则该方程应为
\[
x^4-4x^3+6x^2-4x+1=0,
\]
即 $(x-1)^4=0$. 这个方程的根为 $1,1,1,1$，因此方程组的解为
\[
x_1=1,\quad x_2=1,\quad x_3=1,\quad x_4=1.
\]
\hfill $\square$

\noindent\textbf{例 5.64}\quad 设 $1\le k\le n$，求证：
\[
\sigma_k=\frac1{k!}
\left|\begin{array}{ccccc}
s_1 & 1 & 0 & \cdots & 0\\
s_2 & s_1 & 2 & \cdots & 0\\
\vdots & \vdots & \vdots &        & \vdots\\
s_{k-1} & s_{k-2} & s_{k-3} & \cdots & k-1\\
s_k & s_{k-1} & s_{k-2} & \cdots & s_1
\end{array}\right|,
\]
\[
s_k=
\left|\begin{array}{ccccc}
\sigma_1 & 1 & 0 & \cdots & 0\\
2\sigma_2 & \sigma_1 & 1 & \cdots & 0\\
\vdots & \vdots & \vdots &        & \vdots\\
(k-1)\sigma_{k-1} & \sigma_{k-2} & \sigma_{k-3} & \cdots & 1\\
k\sigma_k & \sigma_{k-1} & \sigma_{k-2} & \cdots & \sigma_1
\end{array}\right|.
\]

\noindent\textbf{证法 1}\quad 由 Newton 公式可得下列线性方程组（将 $\sigma_i$ 看成是未知数）：
\[
\left\{
\begin{array}{l}
\sigma_1=s_1,\\
s_1\sigma_1-2\sigma_2=s_2,\\
s_2\sigma_1-s_1\sigma_2+3\sigma_3=s_3,\\
\cdots\cdots\cdots\\
s_{k-2}\sigma_1-s_{k-3}\sigma_2+\cdots+(-1)^{k-2}(k-1)\sigma_{k-1}=s_{k-1},\\
s_{k-1}\sigma_1-s_{k-2}\sigma_2+\cdots+(-1)^{k-1}k\sigma_k=s_k.
\end{array}
\right.
\]
该方程组的系数行列式为
\[
|A|=
\left|\begin{array}{ccccc}
1 & 0 & 0 & \cdots & 0\\
s_1 & -2 & 0 & \cdots & 0\\
s_2 & -s_1 & 3 & \cdots & 0\\
\vdots & \vdots & \vdots &        & \vdots\\
s_{k-2} & -s_{k-3} & s_{k-4} & \cdots & 0\\
s_{k-1} & -s_{k-2} & s_{k-3} & \cdots & (-1)^{k-1}k
\end{array}\right|
=(-1)^{\frac12 k(k-1)}k!.
\]
又
\[
|A_k|=
\left|\begin{array}{ccccc}
1 & 0 & 0 & \cdots & s_1\\
s_1 & -2 & 0 & \cdots & s_2\\
s_2 & -s_1 & 3 & \cdots & s_3\\
\vdots & \vdots & \vdots &        & \vdots\\
s_{k-2} & -s_{k-3} & s_{k-4} & \cdots & s_{k-1}\\
s_{k-1} & -s_{k-2} & s_{k-3} & \cdots & s_k
\end{array}\right|.
\]
将 $|A_k|$ 的最后一列经 $k-1$ 次相邻对换后换到第一列，再用 $(-1)^{i-2}$ 依次乘以第 $i(3\le i\le k)$ 列，得到
\[
|A_k|=(-1)^{k-1+\frac12(k-1)(k-2)}
\left|\begin{array}{ccccc}
s_1 & 1 & 0 & 0 & \cdots\\
s_2 & s_1 & 2 & 0 & \cdots\\
s_3 & s_2 & s_1 & 3 & \cdots\\
\vdots & \vdots & \vdots & \vdots & \vdots\\
s_{k-1} & s_{k-2} & s_{k-3} & s_{k-4} & \cdots\ k-1\\
s_k & s_{k-1} & s_{k-2} & s_{k-3} & \cdots\ s_1
\end{array}\right|.
\]
于是由 Cramer 法则可得
\[
\sigma_k=\frac{|A_k|}{|A|}=\frac1{k!}
\left|\begin{array}{ccccc}
s_1 & 1 & 0 & \cdots & 0\\
s_2 & s_1 & 2 & \cdots & 0\\
\vdots & \vdots & \vdots &        & \vdots\\
s_{k-1} & s_{k-2} & s_{k-3} & \cdots & k-1\\
s_k & s_{k-1} & s_{k-2} & \cdots & s_1
\end{array}\right|.
\]
另一结论类似可得.

\noindent\textbf{证法 2}\quad 将 $\sigma_k$ 的表达式中的行列式的第 $i(2\le i\le k)$ 列乘以 $(-1)^{i-1}\sigma_{i-1}$ 都加到第一列上，由 Newton 公式可知，所得行列式的第一列除最后一项外都为零，再按第一列展开就可得到第一个结论. 将 $s_k$ 的表达式中的行列式的第 $i(2\le i\le k)$ 列乘以 $(-1)^{i-1}s_{i-1}$ 都加到第一列上，由 Newton 公式可知，所得行列式的第一列除最后一项外都为零，再按第一列展开就可得到第二个结论. \hfill $\square$

\subsection*{§5.10\quad 结式与判别式}

求结式和判别式通常要计算行列式，一般来说计算它们并不容易. 但在一些比较简单的情形，用结式的定义往往就能得到确定的结果. 下面是计算结式和判别式的几个例子.

\noindent\textbf{例 5.65}\quad 求 $n$ 次多项式 $x^n+px+q(n>1)$ 的判别式.

\noindent\textbf{解}\quad 设 $f(x)=x^n+px+q$，$f'(x)=nx^{n-1}+p$. 计算下列 $2n-1$ 阶行列式：
\[
R(f,f')=
\left|\begin{array}{ccccccc}
1 & 0 & \cdots & 0 & p & q & \\
 & 1 & 0 & \cdots & 0 & p & q\\
 & & \ddots & \ddots & & \ddots & \ddots\\
 & & & 1 & 0 & \cdots & 0\quad p\quad q\\
n & 0 & \cdots & 0 & p & & \\
 & n & 0 & \cdots & 0 & p & \\
 & & n & 0 & \cdots & 0 & p\\
 & & & \ddots & \ddots & \ddots & \\
 & & & & n & 0 & \cdots\quad 0\quad p
\end{array}\right|,
\]
求得
\[
R(f,f')=(1-n)^{n-1}p^n+n^nq^{n-1}.
\]
因此
\[
\Delta(f)=(-1)^{\frac12 n(n-1)}\left((1-n)^{n-1}p^n+n^nq^{n-1}\right).
\]
\hfill $\square$

\noindent\textbf{例 5.66}\quad 求下列多项式的判别式：

(1) $f(x)=x^n+2x+1$;

(2) $f(x)=x^n+2$;

(3) $f(x)=x^{n-1}+x^{n-2}+\cdots+x+1$.

\noindent\textbf{解}\quad (1) 在例 5.65 中令 $p=2,q=1$，可得
\[
\Delta(f)=(-1)^{\frac12 n(n-1)}\left(2^n(1-n)^{n-1}+n^n\right).
\]

(2) 在例 5.65 中令 $p=0,q=2$，可得
\[
\Delta(f)=(-1)^{\frac12 n(n-1)}2^{n-1}n^n.
\]

(3) 在例 5.65 中令 $p=0,q=-1$，可得
\[
\Delta(x^n-1)=(-1)^{\frac12(n-1)(n-2)}n^n.
\]
另一方面，由例 5.73 可得
\[
\Delta(x^n-1)=\Delta((x-1)f(x))=f(1)^2\Delta(f(x)),
\]
因此
\[
\Delta(f(x))=(-1)^{\frac12(n-1)(n-2)}n^{n-2}.
\]
\hfill $\square$

\noindent\textbf{例 5.67}\quad 求证：多项式 $f(x)=x^4+px+q$ 有重因式的充要条件是 $27p^4=256q^3$.

\noindent\textbf{证明}\quad 在例 5.65 中令 $n=4$，可得 $\Delta(f)=-27p^4+256q^3$. 因此 $x^4+px+q$ 有重因式的充要条件是 $27p^4=256q^3$. \hfill $\square$

下面是结式与判别式的一些基本性质，它们的证明并不复杂，只需用定义直接验证即可.

\noindent\textbf{例 5.68}\quad 设 $f(x)=a_0x^n+a_1x^{n-1}+\cdots+a_n$，$g(x)=b_0x^m+b_1x^{m-1}+\cdots+b_m$，其中 $a_0\ne0$，$b_0\ne0$，求证：

(1) $R(f,g)=(-1)^{mn}R(g,f)$;

(2) 若 $a,b$ 为非零常数，$R(af,bg)=a^mb^nR(f,g)$.

\noindent\textbf{证明}\quad (1) 对结式定义中的行列式进行 $mn$ 次行对换，就可以将 $R(f,g)$ 变成 $R(g,f)$. 因此 $R(f,g)=(-1)^{mn}R(g,f)$.

(2) 用结式的行列式定义以及行列式性质即得. \hfill $\square$

\noindent\textbf{例 5.69}\quad 求证：$R(f,g_1g_2)=R(f,g_1)R(f,g_2)$.

\noindent\textbf{证明}\quad 设 $f(x)=a_0x^n+a_1x^{n-1}+\cdots+a_{n-1}x+a_n$，其 $n$ 个根为 $x_1,x_2,\cdots,x_n$. 又设 $g_1(x)$ 和 $g_2(x)$ 分别是 $m_1,m_2$ 次多项式，则
\[
\begin{aligned}
R(f,g_1)&=a_0^{m_1}g_1(x_1)g_1(x_2)\cdots g_1(x_n),\\
R(f,g_2)&=a_0^{m_2}g_2(x_1)g_2(x_2)\cdots g_2(x_n),\\
R(f,g_1g_2)&=a_0^{m_1+m_2}g_1(x_1)g_2(x_1)g_1(x_2)g_2(x_2)\cdots g_1(x_n)g_2(x_n).
\end{aligned}
\]
比较上面 3 个等式即知结论成立. \hfill $\square$

下面的例题说明了判别式的用途.

\noindent\textbf{例 5.70}\quad 设 $f(x)$ 是实系数多项式，求证：

(1) 若 $\Delta(f)<0$，则 $f(x)$ 无重根且有奇数对虚根;

(2) 若 $\Delta(f)>0$，则 $f(x)$ 无重根且有偶数对虚根.

\noindent\textbf{证明}\quad 不妨设 $f(x)$ 为 $n$ 次首一多项式，则
\[
\Delta(f)=\prod_{1\le i<j\le n}(x_i-x_j)^2.
\]
显然若 $\Delta(f)\ne0$，则多项式无重根. 现设 $f(x)$ 的根如下：
\[
x_1,\overline{x_1},\cdots,x_k,\overline{x_k};x_{2k+1},\cdots,x_n,
\]
其中 $x_i(1\le i\le k)$ 为虚根，$\overline{x_i}$ 是 $x_i$ 的共轭虚根，$x_j(2k+1\le j\le n)$ 为实根. 对判别式中各项的情况讨论如下：

若 $i,j\ge2k+1$，$i\ne j$，$x_i,x_j$ 都是实数，因此 $(x_i-x_j)^2>0$.

若 $i\le k$，$j\ge2k+1$，则 $(x_i-x_j)(\overline{x_i}-x_j)$ 是实数，因此 $(x_i-x_j)^2(\overline{x_i}-x_j)^2>0$.

若 $i,j\le k$，$i\ne j$，则 $(x_i-x_j)(x_i-\overline{x_j})(\overline{x_i}-x_j)(\overline{x_i}-\overline{x_j})$ 是实数，因此 $(x_i-x_j)^2(x_i-\overline{x_j})^2(\overline{x_i}-x_j)^2(\overline{x_i}-\overline{x_j})^2>0$.

若 $i\le k$，$x_i-\overline{x_i}$ 为纯虚数，故 $(x_i-\overline{x_i})^2<0$.

因此若 $\Delta(f)>0$，则虚根对数为偶数；若 $\Delta(f)<0$，则虚根对数为奇数. \hfill $\square$

下面几道例题涉及判别式的进一步性质，虽然证明比较繁琐，但思路并不复杂，只需利用判别式的定义公式进行计算即可. 计算的时候请注意根的分类.

\noindent\textbf{例 5.71}\quad 设 $f(x)$ 是数域 $\mathbb{F}$ 上的多项式，已知 $\Delta(f(x))$，求 $\Delta(f(x^2))$.

\noindent\textbf{解法 1}\quad 设 $f(x)=a_0x^n+a_1x^{n-1}+\cdots+a_{n-1}x+a_n$，则由例 5.74 可得
\[
\Delta(f(x^2))=(-4)^n a_0a_n(\Delta(f(x)))^2.
\]

\noindent\textbf{解法 2}\quad 设 $f(x)$ 的根为 $x_1,x_2,\cdots,x_n$，则 $f(x^2)$ 的根为 $\pm\sqrt{x_i}(1\le i\le n)$.
\[
\begin{aligned}
\Delta(f(x^2))
&=a_0^{4n-2}\prod_{1\le i<j\le n}(\sqrt{x_i}+\sqrt{x_j})^2(\sqrt{x_i}-\sqrt{x_j})^2(-\sqrt{x_i}+\sqrt{x_j})^2(-\sqrt{x_i}-\sqrt{x_j})^2\\
&\quad\cdot \prod_{i=1}^n(\sqrt{x_i}+\sqrt{x_i})^2\\
&=a_0^{4n-2}\left(\prod_{1\le i<j\le n}(x_i-x_j)^4\right)\cdot 4^nx_1x_2\cdots x_n\\
&=a_0^{4n-2}\left(\prod_{1\le i<j\le n}(x_i-x_j)^4\right)\cdot 4^n(-1)^n\frac{a_n}{a_0}\\
&=(-4)^na_0a_n(\Delta(f(x)))^2.
\end{aligned}
\]
\hfill $\square$

\noindent\textbf{例 5.72}\quad 设 $f(x)$ 和 $g(x)$ 是次数大于 1 的多项式，求证：
\[
\Delta(f(x)g(x))=\Delta(f(x))\Delta(g(x))R(f,g)^2.
\]

\noindent\textbf{证明}\quad 设
\[
f(x)=a_0x^n+a_1x^{n-1}+\cdots+a_{n-1}x+a_n,\quad
g(x)=b_0x^m+b_1x^{m-1}+\cdots+b_{m-1}x+b_m,
\]
且 $f(x),g(x)$ 的根分别是
\[
x_1,x_2,\cdots,x_n;\ x_{n+1},x_{n+2},\cdots,x_{n+m},
\]
则
\[
\Delta(f(x)g(x))=(a_0b_0)^{2(n+m)-2}\prod_{1\le i<j\le n+m}(x_i-x_j)^2.
\]
现将乘积中的因式作如下分类：若 $i\le n,j\le n$，则
\[
a_0^{2n-2}\prod_{1\le i<j\le n}(x_i-x_j)^2=\Delta(f(x));
\]
若 $i>n,j>n$，则
\[
b_0^{2m-2}\prod_{n+1\le i<j\le n+m}(x_i-x_j)^2=\Delta(g(x));
\]
若 $i\le n,j>n$，则
\[
a_0^{2m}b_0^{2n}\prod_{1\le i\le n<j\le n+m}(x_i-x_j)^2=R(f,g)^2.
\]
因此便有
\[
\Delta(f(x)g(x))=\Delta(f(x))\Delta(g(x))R(f,g)^2.
\]
\hfill $\square$

\noindent\textbf{例 5.73}\quad 设 $g(x)$ 是次数大于 1 的多项式，求证：
\[
\Delta((x-a)g(x))=g(a)^2\Delta(g(x)).
\]

\noindent\textbf{证明}\quad 设 $g(x)=b_0x^m+b_1x^{m-1}+\cdots+b_{m-1}x+b_m$，其根为 $x_1,\cdots,x_m$，则
\[
\begin{aligned}
\Delta((x-a)g(x))
&=b_0^{2m}\prod_{i=1}^m(a-x_i)^2\prod_{1\le i<j\le m}(x_i-x_j)^2\\
&=\left(b_0(a-x_1)\cdots(a-x_m)\right)^2\cdot b_0^{2m-2}\prod_{1\le i<j\le m}(x_i-x_j)^2\\
&=g(a)^2\Delta(g(x)).
\end{aligned}
\]
\hfill $\square$

\noindent\textbf{注}\quad 只有当多项式 $f(x),g(x)$ 的次数都大于 0 时，其结式 $R(f(x),g(x))$ 的定义才有意义. 同理，只有当 $f(x)$ 的次数大于 1 时，其判别式 $\Delta(f(x))$ 的定义才有意义. 当然我们也可以作一些人为的规定，例如，若 $g(x)=c$ 是一个非零常数多项式，则约定 $R(f(x),g(x))=c^n$，其中 $n=\deg f(x)$；若 $f(x)$ 是一个一次多项式，则约定 $\Delta(f(x))=1$. 我们不难发现这些约定可以完美地融入到已证明的关于结式和判别式的结果中. 特别地，例 5.72 和例 5.73 的结论也适合 $\deg f(x)=1$ 或 $\deg g(x)=1$ 的情形，并且例 5.73 也可以看成是例 5.72 的特例. 因此从某种意义上说，这些关于结式和判别式的约定都是自然的.

\noindent\textbf{例 5.74}\quad 设 $f(x)=g(h(x))$，其中 $h(x)$ 是 $m$ 次首一多项式，$g(x)$ 是 $n$ 次首一多项式，其根为 $x_1,x_2,\cdots,x_n$，求证：
\[
\Delta(f(x))=\Delta(g(x))^m\Delta(h(x)-x_1)\Delta(h(x)-x_2)\cdots\Delta(h(x)-x_n).
\]

\noindent\textbf{证法 1}\quad 由假设 $g(x)$ 的根为 $x_1,x_2,\cdots,x_n$，故有
\[
f(x)=g(h(x))=(h(x)-x_1)(h(x)-x_2)\cdots(h(x)-x_n).
\]
又设
\[
h(x)-x_i=(x-u_{i1})(x-u_{i2})\cdots(x-u_{im}),\quad 1\le i\le n,
\]
于是 $u_{ij}(1\le i\le n,1\le j\le m)$ 就是 $f(x)$ 的全部根. 对二重足标引进序如下：$(i,j)<(k,l)$ 当且仅当 $i<k$ 或 $i=k,j<l$. 由题目条件可知 $f(x)$ 是首一多项式，因此
\[
\Delta(f(x))=\prod_{1\le(i,j)<(i',j')\le(n,m)}(u_{ij}-u_{i'j'})^2.
\]
对上式乘积中的因子进行分类. 第一类（第一个足标相同）：
\[
\Delta(h(x)-x_i)=\prod_{1\le j<j'\le m}(u_{ij}-u_{ij'})^2,
\]
因此
\[
\prod_{i=1}^n\Delta(h(x)-x_i)=\prod_{i=1}^n\prod_{1\le j<j'\le m}(u_{ij}-u_{ij'})^2.
\]
第二类（第一个足标不同）：注意到对固定的 $i$，$u_{ij}$ 是 $h(x)-x_i$ 的根，因此 $h(u_{ij})=x_i$.

又
\[
\begin{aligned}
h(u_{i1})-x_{i'}&=(u_{i1}-u_{i'1})(u_{i1}-u_{i'2})\cdots(u_{i1}-u_{i'm}),\\
h(u_{i2})-x_{i'}&=(u_{i2}-u_{i'1})(u_{i2}-u_{i'2})\cdots(u_{i2}-u_{i'm}),\\
&\cdots\cdots\cdots\\
h(u_{im})-x_{i'}&=(u_{im}-u_{i'1})(u_{im}-u_{i'2})\cdots(u_{im}-u_{i'm}),
\end{aligned}
\quad i+1\le i'\le n.
\]
上述诸式之积等于
\[
(h(u_{i1})-x_{i'})(h(u_{i2})-x_{i'})\cdots(h(u_{im})-x_{i'})
=(x_i-x_{i'})^m.
\]
因此
\[
\prod_{\substack{1\le(i,j)<(i',j')\le(n,m)\\ i\ne i'}}(u_{ij}-u_{i'j'})^2
=\prod_{1\le i<i'\le m}(x_i-x_{i'})^{2m}
=\Delta(g(x))^m.
\]
综上所述，便有
\[
\Delta(f(x))=\Delta(g(x))^m\Delta(h(x)-x_1)\Delta(h(x)-x_2)\cdots\Delta(h(x)-x_n).
\]

\noindent\textbf{证法 2}\quad 由例 5.69，我们有
\[
\Delta(f)=(-1)^{\frac12 mn(mn-1)}R(f,f'),
\]
\[
\begin{aligned}
R(f,f')
&=R(g(h(x)),g'(h(x))h'(x))\\
&=R(g(h(x)),g'(h(x)))R(g(h(x)),h'(x)),
\end{aligned}
\]
\[
\begin{aligned}
R(g(h(x)),g'(h(x)))
&=\prod_{i=1}^n\prod_{j=1}^m g'(h(u_{ij}))\\
&=\prod_{i=1}^n g'(x_i)^m
=(-1)^{\frac12 mn(n-1)}\Delta(g(x))^m,
\end{aligned}
\]
\[
\begin{aligned}
R(g(h(x)),h'(x))
&=\prod_{i=1}^n\prod_{j=1}^m h'(u_{ij})\\
&=(-1)^{\frac12 mn(m-1)}\prod_{i=1}^n\Delta(h(x)-x_i).
\end{aligned}
\]
因为
\[
\frac12 mn(mn-1)-\frac12 mn(n-1)-\frac12 mn(m-1)=\frac12 m(m-1)n(n-1)
\]
是一个偶数，因此结论成立. \hfill $\square$

\noindent\textbf{例 5.75}\quad 求参数曲线
\[
\left\{
\begin{array}{l}
x=2(t+1)/(t^2+1),\\
y=t^2/(2t-1)
\end{array}
\right.
\]
的直角坐标方程.

\noindent\textbf{解}\quad 去分母得方程组
\[
\left\{
\begin{array}{l}
xt^2-2t+(x-2)=0,\\
t^2-2yt+y=0.
\end{array}
\right.
\]
令
\[
\left\{
\begin{array}{l}
f(t)=xt^2-2t+(x-2),\\
g(t)=t^2-2yt+y,
\end{array}
\right.
\]
由 $t$ 决定的参数曲线上的一点相当于方程组有公共根，因此
\[
R(f,g)=
\left|\begin{array}{cccc}
x & -2 & x-2 & 0\\
0 & x & -2 & x-2\\
1 & -2y & y & 0\\
0 & 1 & -2y & y
\end{array}\right|=0.
\]
求出行列式可得该曲线的直角坐标方程为
\[
5x^2y^2-2x^2y-12xy^2+x^2-4x+12y+4=0.
\]
\hfill $\square$

\subsection*{§5.11\quad 互素多项式的应用}

在高等代数的框架中，多项式理论起到了一个承上启下的作用. 一方面，多项式理论是即将阐述的相似标准型理论的基石；另一方面，它也联系起了前面阐述的矩阵理论和线性空间理论. 下面将通过几个典型例题来看一看互素多项式的相关应用.

\noindent\textbf{例 5.76}\quad 设 $f(x),g(x)$ 是数域 $\mathbb{K}$ 上的互素多项式，$A$ 是 $\mathbb{K}$ 上的 $n$ 阶方阵，满足 $f(A)=O$，证明：$g(A)$ 是可逆矩阵.

\noindent\textbf{证明}\quad 根据假设，存在 $\mathbb{K}$ 上的多项式 $u(x),v(x)$，使得
\[
f(x)u(x)+g(x)v(x)=1.
\]
在上式中代入 $x=A$，可得恒等式
\[
f(A)u(A)+g(A)v(A)=I_n.
\]
因为 $f(A)=O$，故有 $g(A)v(A)=I_n$，从而 $g(A)$ 是非异阵且 $g(A)^{-1}=v(A)$. \hfill $\square$

利用例 5.76 可以证明一大类可逆矩阵的问题，比如例 2.20，而下面的例题则是例 3.70 和例 3.71 的推广.

\noindent\textbf{例 5.77}\quad 设 $f(x),g(x)$ 是数域 $\mathbb{K}$ 上的互素多项式，$A$ 是 $\mathbb{K}$ 上的 $n$ 阶方阵，证明：$f(A)g(A)=O$ 的充要条件是 $r(f(A))+r(g(A))=n$.

\noindent\textbf{证明}\quad 根据假设，存在 $\mathbb{K}$ 上的多项式 $u(x),v(x)$，使得
\[
f(x)u(x)+g(x)v(x)=1.
\]
在上式中代入 $x=A$，可得恒等式
\[
f(A)u(A)+g(A)v(A)=I_n.
\]
考虑如下分块矩阵的初等变换：
\[
\left(\begin{array}{cc}
f(A) & O\\
O & g(A)
\end{array}\right)
\to
\left(\begin{array}{cc}
f(A) & f(A)u(A)\\
O & g(A)
\end{array}\right)
\to
\left(\begin{array}{cc}
f(A) & I_n\\
O & g(A)
\end{array}\right)
\to
\]
\[
\left(\begin{array}{cc}
f(A) & I_n\\
-f(A)g(A) & O
\end{array}\right)
\to
\left(\begin{array}{cc}
O & I_n\\
-f(A)g(A) & O
\end{array}\right)
\to
\left(\begin{array}{cc}
f(A)g(A) & O\\
O & I_n
\end{array}\right),
\]
故有 $r(f(A))+r(g(A))=r(f(A)g(A))+n$，从而结论得证. \hfill $\square$

例 5.78 告诉我们：多项式的互素因式分解可以诱导出空间的直和分解，从几何层面上看，这就是相似标准型理论原始的出发点. 另外，例 5.78 也是例 4.54 和第 4 章解答题 9 的推广.

\noindent\textbf{例 5.78}\quad 设 $f(x),g(x)$ 是数域 $\mathbb{K}$ 上的互素多项式，$\varphi$ 是 $\mathbb{K}$ 上 $n$ 维线性空间 $V$ 上的线性变换，满足 $f(\varphi)g(\varphi)=0$，证明：$V=V_1\oplus V_2$，其中 $V_1=\ker f(\varphi)$，$V_2=\ker g(\varphi)$.

\noindent\textbf{证明}\quad 根据假设，存在 $\mathbb{K}$ 上的多项式 $u(x),v(x)$，使得
\[
f(x)u(x)+g(x)v(x)=1.
\]
在上式中代入 $x=\varphi$，可得恒等式
\[
f(\varphi)u(\varphi)+g(\varphi)v(\varphi)=I_V.
\]
对任意的 $\alpha\in V$，由上式可得
\[
\alpha=f(\varphi)u(\varphi)(\alpha)+g(\varphi)v(\varphi)(\alpha),
\]
注意到 $f(\varphi)u(\varphi)(\alpha)\in\ker g(\varphi)$，$g(\varphi)v(\varphi)(\alpha)\in\ker f(\varphi)$，故有 $V=V_1+V_2$. 任取 $\beta\in V_1\cap V_2$，由上式可得
\[
\beta=u(\varphi)f(\varphi)(\beta)+v(\varphi)g(\varphi)(\beta)=0,
\]
故有 $V_1\cap V_2=0$，因此 $V=V_1\oplus V_2$. \hfill $\square$

下面的例题不仅把例 3.31 推广到更一般的情形，而且利用互素多项式的性质极大地简化了原来的证明，使读者能看清问题的本质.

\noindent\textbf{例 5.79}\quad 设 $\mathbb{Q}(\sqrt[n]{2})=\{a_0+a_1\sqrt[n]{2}+a_2\sqrt[n]{4}+\cdots+a_{n-1}\sqrt[n]{2^{n-1}}\mid a_i\in\mathbb{Q},0\le i\le n-1\}$，证明：$\mathbb{Q}(\sqrt[n]{2})$ 是一个数域，并求 $\mathbb{Q}(\sqrt[n]{2})$ 作为 $\mathbb{Q}$ 上线性空间的一组基.

\noindent\textbf{证明}\quad 设 $f(x)=x^n-2$，由 Eisenstein 判别法可知 $f(x)$ 在 $\mathbb{Q}$ 上不可约，从而 $f(x)$ 是 $\sqrt[n]{2}$ 的极小多项式. 我们先证明：$a_0+a_1\sqrt[n]{2}+\cdots+a_{n-1}\sqrt[n]{2^{n-1}}=0$ 的充要条件是 $a_0=a_1=\cdots=a_{n-1}=0$. 充分性是显然的，现证必要性. 令 $g(x)=a_0+a_1x+\cdots+a_{n-1}x^{n-1}$，则 $g(\sqrt[n]{2})=0$，由极小多项式的基本性质可得 $f(x)\mid g(x)$. 因为 $g(x)$ 的次数小于 $n$，故只能是 $g(x)=0$，即 $a_0=a_1=\cdots=a_{n-1}=0$.

利用 $(\sqrt[n]{2})^n=2$ 容易验证，$\mathbb{Q}(\sqrt[n]{2})$ 中任意两个数的加法、减法和乘法都是封闭的. 要证明 $\mathbb{Q}(\sqrt[n]{2})$ 是数域，只要证明除法或者取倒数封闭即可. 任取 $\mathbb{Q}(\sqrt[n]{2})$ 中的非零数 $\alpha=a_0+a_1\sqrt[n]{2}+\cdots+a_{n-1}\sqrt[n]{2^{n-1}}\ne0$，由上面的讨论可知 $a_0,a_1,\cdots,a_{n-1}$ 不全为零. 令 $g(x)=a_0+a_1x+\cdots+a_{n-1}x^{n-1}$，则 $\alpha=g(\sqrt[n]{2})$. 因为 $f(x)$ 不可约且 $g(x)\ne0$ 的次数小于 $n$，故 $f(x)$ 与 $g(x)$ 互素，由例 5.6 可知，存在有理系数多项式 $u(x),v(x)$，使得
\[
f(x)u(x)+g(x)v(x)=1,\quad \deg v(x)<\deg f(x)=n.
\]
在上式中代入 $x=\sqrt[n]{2}$，可得 $g(\sqrt[n]{2})v(\sqrt[n]{2})=1$，于是 $\alpha^{-1}=v(\sqrt[n]{2})\in\mathbb{Q}(\sqrt[n]{2})$. 因此，$\mathbb{Q}(\sqrt[n]{2})$ 是数域.

由 $\mathbb{Q}(\sqrt[n]{2})$ 的定义可知，$\mathbb{Q}(\sqrt[n]{2})$ 中任一元都是 $1,\sqrt[n]{2},\cdots,\sqrt[n]{2^{n-1}}$ 的 $\mathbb{Q}$-线性组合；又由开始的讨论可知，$1,\sqrt[n]{2},\cdots,\sqrt[n]{2^{n-1}}$ 是 $\mathbb{Q}$-线性无关的，因此它们构成了 $\mathbb{Q}(\sqrt[n]{2})$ 作为 $\mathbb{Q}$ 上线性空间的一组基. 特别地，$\dim_{\mathbb{Q}}\mathbb{Q}(\sqrt[n]{2})=n$. \hfill $\square$

最后，我们给出一道线性变换的例题，它是第 4 章解答题 12 的推广. 我们发现利用互素多项式的性质来证明，不仅避开了繁琐的计算，而且还得到了今后各章节中处理此类问题的一个共同方法.

\noindent\textbf{例 5.80}\quad 设 $f(x)=x^n+a_1x^{n-1}+\cdots+a_{n-1}x+a_n$ 是数域 $\mathbb{K}$ 上的不可约多项式，$\varphi$ 是 $\mathbb{K}$ 上 $n$ 维线性空间 $V$ 上的线性变换，$\alpha_1\ne0,\alpha_2,\cdots,\alpha_n$ 是 $V$ 中的向量，满足
\[
\varphi(\alpha_1)=\alpha_2,\ \varphi(\alpha_2)=\alpha_3,\ \cdots,\ \varphi(\alpha_{n-1})=\alpha_n,\ \varphi(\alpha_n)=-a_n\alpha_1-a_{n-1}\alpha_2-\cdots-a_1\alpha_n.
\]
证明：$\{\alpha_1,\alpha_2,\cdots,\alpha_n\}$ 是 $V$ 的一组基.

\noindent\textbf{证明}\quad 我们只要证明 $\alpha_1,\alpha_2,\cdots,\alpha_n$ 线性无关即可. 用反证法，设存在不全为零的 $n$ 个数 $c_1,c_2,\cdots,c_n$，使得
\[
c_1\alpha_1+c_2\alpha_2+\cdots+c_n\alpha_n=0,
\]
则有 $(c_1I_V+c_2\varphi+\cdots+c_n\varphi^{n-1})(\alpha_1)=0$. 令 $g(x)=c_1+c_2x+\cdots+c_nx^{n-1}$，则 $g(x)\ne0$ 且 $g(\varphi)(\alpha_1)=0$. 另一方面，由假设容易验证 $f(\varphi)(\alpha_1)=0$. 因为 $f(x)$ 不可约且 $g(x)\ne0$ 的次数小于 $n$，故 $f(x)$ 与 $g(x)$ 互素，从而存在 $\mathbb{K}$ 上的多项式 $u(x),v(x)$，使得
\[
f(x)u(x)+g(x)v(x)=1.
\]
在上式中代入 $x=\varphi$，可得恒等式
\[
f(\varphi)u(\varphi)+g(\varphi)v(\varphi)=I_V,
\]
上式两边同时作用 $\alpha_1$ 可得
\[
\alpha_1=u(\varphi)f(\varphi)(\alpha_1)+v(\varphi)g(\varphi)(\alpha_1)=0,
\]
这与假设 $\alpha_1\ne0$ 矛盾，从而结论得证. \hfill $\square$

\subsection*{§5.12\quad 基础训练}

\subsubsection*{5.12.1\quad 训练题}

\noindent\textbf{一、单选题}

1. 若多项式 $f_1(x),f_2(x),f_3(x)$ 互素，则（\quad ）.

(A) $f_1(x),f_2(x)$ 必互素

(B) $f_1(x),f_2(x)$ 互素，或 $f_1(x),f_3(x)$ 互素，或 $f_2(x),f_3(x)$ 互素

(C) 若 $(f_1(x),f_2(x))=d_1(x),(f_2(x),f_3(x))=d_2(x)$，则 $(d_1(x),d_2(x))=1$

(D) 存在 $u(x),v(x)$，使 $f_3(x)=f_1(x)u(x)+f_2(x)v(x)$

2. $m,n$ 是大于 1 的整数，则 $x^{3m}+x^{3n}$ 除以 $x^2+x+1$ 后的余式为（\quad ）.

(A) $x+1$\quad (B) $0$\quad (C) $1$\quad (D) $2$

3. 若 $(f(x),g(x))=d(x)$，则下面的等式成立的是（\quad ）.

(A) $\left(\dfrac{f(x)}{d(x)},g(x)\right)=1$

(B) $\left(\dfrac{f(x)}{d(x)},\dfrac{g(x)}{d(x)}\right)=1$

(C) $(f(x),d(x))=1$

(D) $(f(x),d(x),g(x))=1$

4. 若 $(f(x),g(x))=1,(f(x),h(x))=1$，则下列多项式中不一定互素的是（\quad ）.

(A) $f(x),f(x)+g(x)$

(B) $f(x),h(x)+g(x)$

(C) $f(x),h(x)g(x)$

(D) $f(x)g(x),f(x)+g(x)$

5. 设 $f(x)$ 是数域 $F$ 上的多项式，又 $K$ 是包含 $F$ 的数域，则（\quad ）.

(A) 若 $f(x)$ 在 $F$ 上不可约，则 $f(x)$ 在 $K$ 上也不可约

(B) 若 $f(x)$ 在 $K$ 上可约，则 $f(x)$ 在 $F$ 上也可约

(C) 若 $f(x)$ 在 $F$ 上有重因式，则 $f(x)$ 在 $K$ 上必有重根

(D) 若 $f(x)$ 在 $K$ 上不可约，则 $f(x)$ 在 $F$ 上也不可约

6. 设 $f(x)$ 是整系数多项式，则下列命题正确的是（\quad ）.

(A) $f(x)$ 有有理根的充要条件是 $f(x)$ 在有理数域上可约

(B) 若既约分数 $\dfrac{p}{q}$ 是 $f(x)$ 的根，则 $p$ 可整除 $f(x)$ 的常数项

(C) 若 $p$ 是素数且能整除 $f(x)$ 的除首项外的所有项系数，则 $f(x)$ 在有理数域上不可约

(D) 若 $f(x)$ 有重因式，则它在有理数域上必有重根

7. 设有理系数多项式 $f(x)=cp_1(x)p_2(x)\cdots p_k(x)$，其中 $c\ne0$，$p_i(x)$ 为互不相同的首一不可约有理系数多项式，则 $f(x)$ 在复数域内（\quad ）.

(A) 无重根\quad (B) 可能有重根\quad (C) 无实根\quad (D) 有 $k$ 个实根

8. 下列多项式在复数域中有重根的是（\quad ）.

(A) $x^n+1$

(B) $x^n+x^{n-1}+\cdots+x+1$

(C) $1+x+\dfrac{x^2}{2!}+\cdots+\dfrac{x^n}{n!}$

(D) $nx^{n+1}-(n+1)x^n+1$

9. 下列多项式在有理数域上不可约的是（\quad ）.

(A) $x^{2n+1}+1$

(B) $x^4+4$

(C) $x^6+x^3+1$

(D) $x^4-3x^3+5x^2+2x-5$

10. 下列命题正确的是（\quad ）.

(A) 若复数 $c$ 是多项式 $f(x)$ 的 $k$ 重根，则 $c$ 是 $f'(x)$ 的 $k-1$ 重根

(B) 若复数 $c$ 是多项式 $f(x)$ 的导数 $f'(x)$ 的 $k$ 重根，则 $c$ 是 $f(x)$ 的 $k+1$ 重根

(C) 若复数 $c$ 是多项式 $f(x)$ 的 $k$ 重根，则 $c$ 也是 $f'(x)$ 的 $k$ 重根

(D) 若 $f(x),f'(x)$ 的最大公因式是 $k$ 次多项式，则 $f(x)$ 有 $k$ 重根

11. 若 $f(x)$ 是有理数域上的可约多项式，则正确的结论应该是（\quad ）.

(A) 由 $f(x)\mid g(x)h(x)$ 可推出 $f(x)\mid g(x)$ 或 $f(x)\mid h(x)$

(B) $f(x)$ 必有有理根

(C) 若 $p_1(x),p_2(x)$ 是 $f(x)$ 的不可约因式，且 $p_1(x)$ 和 $p_2(x)$ 在有理数域内互素，则 $p_1(x),p_2(x)$ 在复数域内无公根

(D) $f(x)$ 的不可约因式的次数不超过 2

12. 设 $f(x)$ 是实数域上的多项式，则错误的结论应该是（\quad ）.

(A) 若 $f(x)$ 的次数是奇数，则它必有实数根

(B) 若 $f(x)$ 可约，则它必有实数根

(C) 若 $f(x)$ 的系数全是正实数，则它没有正实数根

(D) 若 $f(x)$ 可约，则每个不可约因式的次数不超过 2

13. 以 $\sqrt{2}-1+i$ 为根的次数最小的有理系数多项式的次数为（\quad ）.

(A) 2\quad (B) 3\quad (C) 4\quad (D) 6

14. 下列多项式是对称多项式的是（\quad ）.

(A) $x_1^3+x_2^3-x_3^3$

(B) $2x_1^3+x_2^3+x_3^3$

(C) $x_1^3+x_2^3+x_3^3+2x_1x_2+2x_1x_3+2x_2x_3$

(D) $x_1^2+x_2^2+x_3^2+x_1x_2$

15. 下列多项式有公根的是（\quad ）.

(A) $x^3-4x^2+5x-2,\ x^2+x+1$

(B) $x^3-4x^2+5x-2,\ x^2+x-2$

(C) $x^4-x^2+3x-8,\ x^3+x^2+x+1$

(D) $x^4-x^2+3x-8,\ x^3-1$

\noindent\textbf{二、填空题}

1. 设 $f(x)=2x^4-3x^3+4x^2+ax+b,\ g(x)=x^2-3x+1$，若 $f(x)$ 除以 $g(x)$ 后余式等于 $25x-5$，则 $a=($\quad $),\ b=($\quad $)$.

2. 当 $a,b,c$ 适合条件（\quad ）时，$(x^2+c)\mid(x^3+ax+b)$.

3. 设 $x_1,x_2,x_3$ 是多项式 $x^3-6x^2+5x-1$ 的 3 个根，则
\[
(x_1-x_2)^2+(x_1-x_3)^2+(x_2-x_3)^2=(\quad).
\]

4. 设 $f(x)$ 是一个三次首一多项式，若 $f(x)$ 除以 $x-1$ 余 1，除以 $x-2$ 余 2，除以 $x-3$ 余 3，则 $f(x)=($\quad $)$.

5. $x^4+2$ 在实数域上是否可约？（\quad ）

6. 设 2 是多项式 $x^4-2x^3+ax^2+bx-8$ 的二重根，则 $a=($\quad $),\ b=($\quad $)$.

7. 设三次方程 $x^3+px^2+qx+r=0\ (r\ne0)$，以此方程根的倒数为根的三次方程为（\quad ）.

8. 方程 $4x^4-7x^2-5x-1=0$ 的有理根为（\quad ）.

9. 设 $p$ 是素数，则 $x^p+px+p$ 和 $x^2+p$ 的最大公因式是（\quad ）.

10. 已知实系数多项式 $x^3+px+q$ 有一个虚根 $3+2i$，则其余两个根为（\quad ）.

11. $f(x)=x^4-3x^3+2x^2-5x+1$ 被 $x-2$ 除后的余数是（\quad ）.

12. 若二元多项式 $f(x_1,x_2),g(x_1,x_2)$ 在一个无穷集上的值相等，问它们是否相同？（\quad ）

13. 设 $\sigma_1=x_1+x_2+x_3,\sigma_2=x_1x_2+x_1x_3+x_2x_3,\sigma_3=x_1x_2x_3$，则 $x_1^3+x_2^3+x_3^3$ 可用 $\sigma_i$ 表示为（\quad ）.

14. 多项式 $x^3-2x^2-5x+6$ 的判别式值等于（\quad ）.

15. 当实数 $t=($\quad $)$ 时，多项式 $x^3+tx-2$ 有重根.

\noindent\textbf{三、解答题}

1. 设 $V$ 是由数域 $F$ 上的多项式全体组成的线性空间，$\{f_1(x),\cdots,f_m(x)\},\{g_1(x),\cdots,g_n(x)\}$ 是 $V$ 中两组向量. 假设这两组向量等价，即它们可互相线性表示，求证：$\{f_i(x)\}$ 的最大公因子等于 $\{g_j(x)\}$ 的最大公因子.

2. 设 $m(x)$ 是 $\{f_1(x),\cdots,f_n(x)\}$ 的公倍式，求证：$m(x)$ 是最小公倍式的充要条件是
\[
\left\{\frac{m(x)}{f_1(x)},\cdots,\frac{m(x)}{f_n(x)}\right\}
\]
是一组互素的多项式.

3. 设 $f(x),g(x)$ 是数域 K 上的互素多项式，$u(x),v(x)$ 也是 K 上的互素多项式且 $u(x)s(x)-v(x)t(x)=1$. 求证：对任意的正整数 $m,n$，多项式 $u(x)f(x)^m+v(x)g(x)^n$ 和 $t(x)f(x)^m+s(x)g(x)^n$ 必互素.

4. 设 $p_1(x),p_2(x)$ 是两个互素多项式，若 $f(x)$ 除以 $p_1(x)$ 的余式为 $r_1(x)$，$f(x)$ 除以 $p_2(x)$ 的余式为 $r_2(x)$. 设 $u(x),v(x)$ 是使 $p_1(x)u(x)+p_2(x)v(x)=1$ 的多项式. 令
\[
g(x)=r_2(x)p_1(x)u(x)+r_1(x)p_2(x)v(x),
\]
证明：$g(x),f(x)$ 关于 $p_1(x)p_2(x)$ 同余.

5. 求使 $1+x^n+x^{2n}+\cdots+x^{mn}$ 能被 $1+x+x^2+\cdots+x^m$ 整除的所有正整数对 $(m,n)$.

6. 证明：$x^n+ax^{n-m}+b\ (b\ne0)$ 不能有重数大于 2 的重根.

7. 求证：1 是下列多项式的 3 重根：
\[
x^{2n+1}-(2n+1)x^{n+1}+(2n+1)x^n-1.
\]

8. 求证：多项式 $f(x)=a_nx^n+a_{n-1}x^{n-1}+\cdots+a_1x+a_0$ 能被 $(x-1)^{k+1}$ 整除的充要条件是
\[
\begin{cases}
a_0+a_1+a_2+\cdots+a_n=0,\\
a_1+2a_2+\cdots+na_n=0,\\
\cdots\cdots\cdots\cdots,\\
a_1+2^ka_2+\cdots+n^ka_n=0.
\end{cases}
\]

9. 设 $f(x)$ 是数域 $F$ 上的多项式且对任意的 $a,b\in F$，总有 $f(a+b)=f(a)+f(b)$，求证：$f(x)=kx$，其中 $k\in F$.

10. 证明：如果 $(x-1)\mid f(x^n)$，则 $(x^n-1)\mid f(x^n)$.

11. 证明：如果 $(x^2+x+1)\mid(f_1(x^3)+xf_2(x^3))$，则 $(x-1)\mid f_1(x)$ 且 $(x-1)\mid f_2(x)$.

12. 设 $V$ 是数域 $F$ 上全体多项式组成的线性空间，$D$ 是 $V$ 上的线性变换，若 $D$ 适合
\[
(1)\ D(x)=1;\qquad (2)\ D(f(x)g(x))=g(x)D(f(x))+f(x)D(g(x)).
\]
求证：$D$ 就是求导变换.

13. 设 $x^3+px^2+qx+r=0$ 的 3 个根为 $x_1,x_2,x_3$，求一个三次方程其根为 $x_1x_2,x_1x_3,x_2x_3$.

14. 设 $f(x)=x^{2n+1}-1$，$f(x)$ 的不等于 1 的根为 $\omega_1,\omega_2,\cdots,\omega_{2n}$，求证：
\[
(1-\omega_1)(1-\omega_2)\cdots(1-\omega_{2n})=2n+1.
\]

15. 设 $f(x)$ 是次数小于 $n$ 的多项式，$\varepsilon=\cos\dfrac{2\pi}{n}+i\sin\dfrac{2\pi}{n}$ 是 1 的 $n$ 次根，求证：
\[
f(0)=\frac{1}{n}\sum_{i=1}^n f(\varepsilon^i).
\]

16. 求证：方程 $x^8+5x^6+4x^4+2x^2+1=0$ 无实数根.

17. 设有实数 $a,b,c$，求证：$a>0,b>0,c>0$ 的充要条件是：
\[
a+b+c>0,\quad ab+ac+bc>0,\quad abc>0.
\]

18. 设 $f(x)$ 是整系数多项式，既约分数 $\dfrac{p}{q}$ 是 $f(x)$ 的根，求证：$f(x)=(qx-p)g(x)$，其中 $g(x)$ 是一个整系数多项式.

19. 设 $f(x)$ 是 $n$ 次有理系数多项式，若 $k>n$，求证：$\sqrt[k]{2}$ 必不是 $f(x)$ 的根.

20. 写出一个次数最小的首一有理系数多项式，使它有下列根：
\[
1+\sqrt{3},\quad 3+\sqrt{2}i.
\]

21. 设奇数次多项式 $f(x)=(x-a_1)(x-a_2)\cdots(x-a_n)+1$，其中 $a_i$ 是互不相同的整数. 求证：$f(x)$ 在有理数域上不可约.

22. 设 $f(x)$ 是整系数多项式，若有 4 个不同的整数 $a_1,\cdots,a_4$，使得 $f(a_i)=5(1\le i\le4)$，求证：$f(x)$ 没有整数根.

23. 解下列方程组：
\[
\begin{cases}
x_1+x_2+\cdots+x_n=n,\\
x_1^2+x_2^2+\cdots+x_n^2=n,\\
\cdots\cdots\cdots\cdots,\\
x_1^n+x_2^n+\cdots+x_n^n=n.
\end{cases}
\]

24. 计算方程
\[
x^n+(a+b)x^{n-1}+(a^2+ab+b^2)x^{n-2}+\cdots+(a^n+a^{n-1}b+\cdots+b^n)=0
\]
根的方幂和 $s_k=x_1^k+x_2^k+\cdots+x_n^k$，其中 $k\le n$.

25. 当 $c$ 是什么数时，多项式 $x^3+cx+1,x^2+cx+1$ 有公根？

\subsubsection*{5.12.2\quad 训练题答案}

\noindent\textbf{一、单选题}

1. 应选择 (C).

2. 应选择 (D). 注意到 $x^2+x+1$ 可整除 $x^{3k}-1$，而 $x^{3m}+x^{3n}=(x^{3m}-1)+(x^{3n}-1)+2$.

3. 应选择 (B). (C) 和 (D) 显然不正确，(A) 也不正确，例如 $f(x)=x^2,g(x)=x$.

4. 应选择 (B). 事实上，若 $h(x)=f(x)-g(x)$，则 $(f(x),h(x))=1,\ h(x)+g(x)=f(x)$.

5. 应选择 (D). 注意 (C) 不正确，因为 $f(x)$ 在 K 上未必有根.

6. 应选择 (B). 注意 (A) 不正确，例如 $f(x)=(x^2-2)(x^2-3)$.

7. 应选择 (A). 由于 $p_i(x)$ 不可约，故 $p_i(x)$ 无重根. 又由已知条件，$p_i(x)$ 两两互素，所以对任意的 $i\ne j$，存在 $u(x),v(x)$，使得 $p_i(x)u(x)+p_j(x)v(x)=1$，于是 $p_i(x)$ 之间无公根. 因此 $f(x)$ 无重根.

8. 应选择 (D). 计算 $f(x)$ 和 $f'(x)$ 的公因式即可.

9. 应选择 (C). 令 $x=y+1$，由 Eisenstein 判别法可知 $x^6+x^3+1$ 在 $\mathbb{Q}$ 上不可约. (A)，(B)，(D) 中的多项式在有理数域上都可进行因式分解.

10. 应选择 (A)，参考例 5.23. 注意 (B) 不正确，例如 $f(x)=x^2+2x$，$-1$ 是导数 $f'(x)=2x+2$ 的 1 重根，但不是 $f(x)$ 的 2 重根.

11. 应选择 (C)，理由与第 7 题相同.

12. 应选择 (B)，例如 $f(x)=(x^2+1)(x^2+2)$.

13. 应选择 (C)，极小多项式为
\[
(x-\sqrt{2}+1-i)(x-\sqrt{2}+1+i)(x+\sqrt{2}+1-i)(x+\sqrt{2}+1+i).
\]

14. 应选择 (C).

15. 应选择 (B). 可以计算结式进行验证，也可以计算低次多项式的根，再代入高次多项式进行验证.

\noindent\textbf{二、填空题}

1. 用带余除法即可求得 $a=-5,b=6$.

2. 同上题求法可得 $a=c,b=0$.

3. 用 Vieta 定理即可求得答案为 42.

4. $f(x)=(x-1)(x-2)(x-3)+x=x^3-6x^2+12x-6$.

5. 因为不可约实系数多项式的次数不超过 2，故此多项式必可约.

6. 将 2 代入多项式及其导数可知，它们都应等于零，可求得 $a=-6,b=16$.

7. $rx^3+qx^2+px+1$.

8. $-\dfrac{1}{2}$.

9. 由 Eisenstein 判别法可知，这两个多项式在有理数域上不可约. 又后者不能整除前者，故两者必互素，因此最大公因式为 1.

10. 实系数方程虚根必成对出现，因此另一个虚根为 $3-2i$，再由 Vieta 定理可得另一根为 $-6$.

11. 用余数定理求得余数为 $-9$.

12. 不一定. 例如，$f(x_1,x_2)=x_1x_2,\ g(x_1,x_2)=x_1^2+x_1x_2$，它们在无穷集 $\{(0,x_2)\mid x_2\in K\}$ 上的值相等.

13. $\sigma_1^3-3\sigma_1\sigma_2+3\sigma_3$.

14. 900.

15. $-3$.

\noindent\textbf{三、解答题}

1. 设 $d_1(x),d_2(x)$ 分别是 $\{f_i(x)\}$ 及 $\{g_j(x)\}$ 的最大公因式，则由已知可推出 $d_1(x)\mid g_j(x)$，故 $d_1(x)\mid d_2(x)$. 同理可证 $d_2(x)\mid d_1(x)$，因此 $d_1(x)=d_2(x)$.

2. 设 $m(x)=f_i(x)g_i(x)$，假设 $d(x)$ 是 $g_i(x)=\dfrac{m(x)}{f_i(x)}$ 的最大公因式，则 $\dfrac{m(x)}{d(x)}$ 仍是 $f_i(x)$ 的公倍式. 如果 $d(x)\ne1$，则 $m(x)$ 不是 $f_i(x)$ 的最小公倍式. 反之，若 $m(x)$ 不是最小公倍式，设 $m(x)=h(x)t(x)$，其中 $h(x)$ 是 $f_i(x)$ 的最小公倍式，则 $t(x)$ 是 $g_i(x)$ 的公因式，于是 $d(x)\ne1$.

3. 令 $\varphi(x)=u(x)f(x)^m+v(x)g(x)^n$ 和 $\psi(x)=t(x)f(x)^m+s(x)g(x)^n$，则
\[
f(x)^m=s(x)\varphi(x)-v(x)\psi(x),\qquad
g(x)^n=-t(x)\varphi(x)+u(x)\psi(x).
\]
假设 $p(x)$ 是 $\varphi(x)$ 和 $\psi(x)$ 的不可约因式，则由上式可得 $p(x)\mid f(x)^m,\ p(x)\mid g(x)^n$，从而 $p(x)\mid f(x),\ p(x)\mid g(x)$，这与假设 $(f(x),g(x))=1$ 相矛盾.

4. 设 $f(x)=p_1(x)q_1(x)+r_1(x),\ f(x)=p_2(x)q_2(x)+r_2(x)$，则
\[
\begin{aligned}
f(x)-g(x)
&=f(x)(p_1(x)u(x)+p_2(x)v(x))-r_2(x)p_1(x)u(x)-r_1(x)p_2(x)v(x)\\
&=(f(x)-r_2(x))p_1(x)u(x)+(f(x)-r_1(x))p_2(x)v(x)\\
&=p_2(x)q_2(x)p_1(x)u(x)+p_1(x)q_1(x)p_2(x)v(x)\\
&=p_1(x)p_2(x)(q_2(x)u(x)+q_1(x)v(x)),
\end{aligned}
\]
故 $p_1(x)p_2(x)\mid(f(x)-g(x))$. 也可利用中国剩余定理的证明过程直接得到结论.

5. 由已知，需要
\[
\frac{x^{m+1}-1}{x-1}\mid\frac{x^{(m+1)n}-1}{x^n-1},
\]
即
\[
\frac{x^{m+1}-1}{x-1}\cdot\frac{x^n-1}{x-1}\mid\frac{x^{(m+1)n}-1}{x-1}. \tag{5.2}
\]
注意到 $(x^{m+1}-1)\mid(x^{(m+1)n}-1)$，$(x^n-1)\mid(x^{(m+1)n}-1)$，故有
\[
\frac{x^{m+1}-1}{x-1}\mid\frac{x^{(m+1)n}-1}{x-1},\qquad
\frac{x^n-1}{x-1}\mid\frac{x^{(m+1)n}-1}{x-1}.
\]
若 $m+1$ 和 $n$ 互素，则由例 5.14 以及互素多项式的性质可知，$\dfrac{x^{m+1}-1}{x-1}$ 和 $\dfrac{x^n-1}{x-1}$ 互素，从而 (5.2) 式成立. 反之，若 (5.2) 式成立，则由 $\dfrac{x^{(m+1)n}-1}{x-1}$ 无重根可知，$\dfrac{x^{m+1}-1}{x-1},\dfrac{x^n-1}{x-1}$ 必互素，从而由互素多项式的性质可得 $(x^{m+1}-1,x^n-1)=x-1$，再由例 5.14 可知，$m+1$ 和 $n$ 互素. 因此所有满足 $m+1$ 和 $n$ 互素的正整数对 $(m,n)$ 即为所求.

6. 该多项式的导数为 $nx^{n-1}+(n-m)ax^{n-m-1}=x^{n-m-1}(nx^m+(n-m)a)$. 和原多项式比较，若 $a=0$，则 $x=0$ 不是原多项式的根，从而原多项式无重根；若 $a\ne0$，则 $x=0$ 也不是原多项式的根，而 $nx^m+(n-m)a=0$ 只有单根，因此原多项式的重根数不大于 2.

7. 容易验证 1 是多项式的根，又是其导数及其二阶导数的根，但不是其三阶导数的根，从而由例 5.23 即得结论.

8. 由例 5.23 可知，$f(x)$ 能被 $(x-1)^{k+1}$ 整除当且仅当 $f(1)=f'(1)=\cdots=f^{(k)}(1)=0$. 定义 $g_0(x)=f(x),\ g_1(x)=f'(x),\ g_i(x)=(xg_{i-1}(x))'(i\ge2)$，则通过简单的验证可知，$f(1)=f'(1)=\cdots=f^{(k)}(1)=0$ 当且仅当 $g_0(1)=g_1(1)=\cdots=g_k(1)=0$，而后者即为题中所给条件.

9. 设 $f(1)=k$，则 $f(n)=kn(n=1,2,\cdots)$. 于是多项式 $f(x)-kx$ 有无穷多根，从而 $f(x)-kx=0$.

10. 由余数定理可得 $f(1)=0$，故 $(x-1)\mid f(x)$，于是 $(x^n-1)\mid f(x^n)$.

11. 设 1 的三次根为 $1,\omega,\omega^2$，则 $f_1(1)+\omega f_2(1)=0,\ f_1(1)+\omega^2f_2(1)=0$. 因此 $f_1(1)=f_2(1)=0$，由余数定理即得结论.

12. 用归纳法容易证明 $D(x^n)=nx^{n-1}$，再利用线性即得结论.

13. 用 Vieta 定理可求得 $f(x)=x^3-qx^2+prx-r^2$.

14. 分解因式 $f(x)=(x-1)(x^{2n}+x^{2n-1}+\cdots+1)=(x-1)(x-\omega_1)(x-\omega_2)\cdots(x-\omega_{2n})$，因此
\[
x^{2n}+x^{2n-1}+\cdots+1=(x-\omega_1)(x-\omega_2)\cdots(x-\omega_{2n}).
\]
令 $x=1$ 代入即得结论.

15. 设 $f(x)=a_0+a_1x+\cdots+a_mx^m(m<n)$，则
\[
\sum_{i=1}^n f(\varepsilon^i)=na_0+a_1\sum_{i=1}^n\varepsilon^i+\cdots+a_m\sum_{i=1}^n\varepsilon^{mi}.
\]
注意到 $\varepsilon^j(1\le j<n)$ 都是 1 的 $n$ 次根，都适合多项式 $x^{n-1}+x^{n-2}+\cdots+1$，于是 $\sum_{i=1}^n\varepsilon^{ji}=0$，由此即得结论.

16. 将任一实数代入方程都得到大于零的数，因此方程无实数根.

17. 只需证明充分性. 设 $a+b+c=p,\ ab+ac+bc=q,\ abc=r$，则 $f(x)=x^3-px^2+qx-r$ 以 $a,b,c$ 为实根. 该方程没有负根，零也不是根.

18. 设 $f(x)=(qx-p)g(x)$，其中 $g(x)$ 是有理系数多项式. 不妨设 $g(x)=ch(x)$，$h(x)$ 是本原多项式，$c$ 是有理数，则 $f(x)=c(qx-p)h(x)$. 由 Gauss 引理可知，$(qx-p)h(x)$ 也是本原多项式，如果 $c$ 是分数，将和 $f(x)$ 是整系数多项式矛盾. 因此 $c$ 是整数，从而 $g(x)$ 是整系数多项式.

19. 由 Eisenstein 判别法可知，$x^k-2$ 在 $\mathbb{Q}$ 上不可约，从而它是 $\sqrt[k]{2}$ 的极小多项式. 若 $\sqrt[k]{2}$ 是 $f(x)$ 的根，则由极小多项式的基本性质可得 $(x^k-2)\mid f(x)$，但 $\deg f(x)=n<k$，矛盾.

20. 容易验证 $1+\sqrt{3}$ 的极小多项式是 $x^2-2x-2,\ 3+\sqrt{2}i$ 的极小多项式是 $x^2-6x+11$. 若有理系数多项式 $f(x)$ 有根 $1+\sqrt{3},3+\sqrt{2}i$，则由极小多项式的基本性质可得 $(x^2-2x-2)\mid f(x),\ (x^2-6x+11)\mid f(x)$. 因为 $(x^2-2x-2,x^2-6x+11)=1$，故有 $(x^2-2x-2)(x^2-6x+11)\mid f(x)$，于是满足条件的次数最小的首一有理系数多项式为
\[
(x^2-2x-2)(x^2-6x+11)=x^4-8x^3+21x^2-10x-22.
\]

21. 用反证法，设 $f(x)=g(x)h(x)$ 是两个次数小于 $n$ 的整系数多项式的乘积，则 $g(a_i)h(a_i)=1$. 因为 $g(x),h(x)$ 是整系数多项式，故必有 $g(a_i)=1,h(a_i)=1$ 或 $g(a_i)=-1,h(a_i)=-1$. 无论怎样，都有 $g(a_i)=h(a_i)(1\le i\le n)$. 由于 $g(x),h(x)$ 的次数小于 $n$，故 $g(x)=h(x)$，于是 $f(x)=g(x)^2$，这和 $n$ 是奇数矛盾.

22. 设 $f(x)=(x-a_1)(x-a_2)(x-a_3)(x-a_4)g(x)+5$，因为 $f(x)-5$ 是整系数多项式，$(x-a_1)(x-a_2)(x-a_3)(x-a_4)$ 是本原多项式，所以 $g(x)$ 必是整系数多项式. 用反证法，假设 $f(x)$ 有整数根 $c$，则 $-5=(c-a_1)(c-a_2)(c-a_3)(c-a_4)g(c)$. 注意到 $-5$ 仅有因子 $\pm1,\pm5$，而 $c-a_i(1\le i\le4)$ 是 4 个不同的整数，从而它们只能分别是 $\pm1,\pm5$，于是 $g(c)=-\dfrac{1}{5}\notin\mathbb{Z}$，矛盾.

23. 类比例 5.63 可得方程的解为 $x_1=x_2=\cdots=x_n=1$.

24. 对 $k$ 用数学归纳法以及 Newton 公式即可证明 $s_k=-(a^k+b^k)$.

25. 利用结式可解得 $c=-2$.

\end{document}
